\documentclass[12pt,a4paper]{report}
\usepackage[utf8]{inputenc}
\usepackage[T1]{fontenc}
\usepackage{geometry}
\usepackage{booktabs}
\usepackage{amsmath}
\usepackage{amssymb}
\usepackage{graphicx}
\usepackage{xcolor}
\usepackage{array}
\usepackage{longtable}
\usepackage{multirow}
\usepackage{hyperref}
\usepackage{fancyhdr}
\usepackage{titlesec}
\usepackage{enumitem}
\usepackage{caption}
\usepackage{subcaption}
\usepackage{float}
\usepackage{tabularx}
\usepackage{colortbl}
\usepackage{pdflscape}
\usepackage{rotating}
\usepackage{listings}
\usepackage{courier}
\usepackage{setspace}
\usepackage{microtype}

\geometry{margin=2.5cm, top=3cm, bottom=3cm}

%% Colours
\definecolor{airblue}{RGB}{30,100,180}
\definecolor{evapgreen}{RGB}{34,139,34}
\definecolor{chilledred}{RGB}{180,30,30}
\definecolor{lightgray}{RGB}{245,245,245}
\definecolor{darkgray}{RGB}{80,80,80}
\definecolor{passgreen}{RGB}{0,128,0}
\definecolor{failred}{RGB}{200,0,0}
\definecolor{warnorng}{RGB}{200,100,0}
\definecolor{headerbg}{RGB}{30,60,100}

%% Header / Footer
\pagestyle{fancy}
\fancyhf{}
\rhead{\small FYP: Data Center Cooling Strategies}
\lhead{\small Aleeza Hussain}
\cfoot{\thepage}
\renewcommand{\headrulewidth}{0.4pt}

%% Hyperref
\hypersetup{colorlinks=true,linkcolor=black,urlcolor=blue,citecolor=black,
  pdftitle={FYP Results Chapter},pdfauthor={Aleeza Hussain}}

%% Section formatting
\titleformat{\chapter}[display]
  {\normalfont\huge\bfseries\color{headerbg}}{\chaptertitlename\ \thechapter}{20pt}{\Huge}
\titleformat{\section}{\normalfont\Large\bfseries\color{headerbg}}{\thesection}{1em}{}
\titleformat{\subsection}{\normalfont\large\bfseries\color{darkgray}}{\thesubsection}{1em}{}
\titleformat{\subsubsection}{\normalfont\normalsize\bfseries}{\thesubsubsection}{1em}{}

%% Convenience macros
\newcommand{\pass}{\textcolor{passgreen}{\textbf{PASS}}}
\newcommand{\fail}{\textcolor{failred}{\textbf{FAIL}}}
\newcommand{\marginal}{\textcolor{warnorng}{\textbf{MARGINAL}}}
\newcommand{\degC}{$^{\circ}$C}
\newcommand{\kWh}{kWh}
\newcommand{\kW}{kW}

\begin{document}
\listoftables
\listoffigures
\newpage

%% ============================================================
\chapter{Reference Values Used Throughout This Report}
\label{ch:reference_values}
%% ============================================================

This chapter consolidates the key numerical values used repeatedly across the Executive Summary,
technique chapters, and comparative analysis sections. All values below are identical to the
values reported in the detailed tables later in this document.

\section{Executive Summary Reference Values}

\begin{table}[H]
\centering
\caption{Reference Values --- Executive Summary Metrics}
\label{tab:ref_exec}
\renewcommand{\arraystretch}{1.3}
\begin{tabular}{lrrrr}
\toprule
\textbf{Metric} & \textbf{Air-Side Econ.} & \textbf{Evaporative} & \textbf{Chilled Water} & \textbf{Best} \\
\midrule
Total Energy (kWh/yr) & 478{,}090 & 141{,}539 & 288{,}475 & Evaporative \\
IT Energy (kWh/yr) & 372{,}125 & 141{,}173 & 248{,}000 & Evaporative \\
Average PUE & 1.285 & 1.002 & 1.163 & Evaporative \\
Average CUE (kgCO$_2$/kWh$_{IT}$) & 0.591 & 0.461 & 0.549 & Evaporative \\
Annual OpEx (USD) & \$97{,}033 & \$20{,}523 & \$42{,}553 & Evaporative \\
Carbon Emissions (kg/yr) & 219{,}921 & 65{,}108 & 158{,}196 & Evaporative \\
Water Usage (L/yr) & 0 & 0 & 515{,}331 & Air-Side / Evap \\
CAPEX (USD) & \$45{,}000 & Low & \$630{,}000 & Evaporative \\
Payback Period (yrs) & 491.9 & N/A & 50.0 & --- \\
ASHRAE Thermal Compliance & Marginal & FAIL & PASS & Chilled Water \\
\bottomrule
\end{tabular}
\end{table}

\section{Air-Side Economizer Reference Values}

\begin{table}[H]
\centering
\caption{Reference Values --- Air-Side Economizer}
\label{tab:ref_air}
\renewcommand{\arraystretch}{1.3}
\begin{tabular}{lrr}
\toprule
\textbf{Metric} & \textbf{Value} & \textbf{Unit} \\
\midrule
Total IT Energy & 372{,}125 & kWh \\
Total Cooling Energy & 105{,}965 & kWh \\
Total Facility Energy & 478{,}090 & kWh \\
Average PUE & 1.285 & --- \\
Average CUE & 0.591 & kgCO$_2$/kWh$_{IT}$ \\
Total Carbon Emissions & 219{,}921 & kg CO$_2$ \\
Water Usage & 0 & Litres \\
CAPEX & \$45{,}000 & USD \\
Annual Electricity Cost & \$69{,}323 & USD \\
Annual Carbon Tax Cost & \$27{,}710 & USD \\
Total Annual OpEx & \$97{,}033 & USD \\
Simple Payback Period & 491.9 & Years \\
\bottomrule
\end{tabular}
\end{table}

\section{Evaporative Cooling Reference Values}

\begin{table}[H]
\centering
\caption{Reference Values --- Evaporative Cooling}
\label{tab:ref_evap}
\renewcommand{\arraystretch}{1.3}
\begin{tabular}{lrr}
\toprule
\textbf{Metric} & \textbf{Value} & \textbf{Unit} \\
\midrule
Total IT Energy & 141{,}173 & kWh \\
Fan Energy (Auxiliary) & 366 & kWh \\
DX Compressor Energy & 0 & kWh \\
Pump Energy & 0 & kWh \\
Total Facility Energy & 141{,}539 & kWh \\
Average PUE & 1.002 & --- \\
Peak PUE & 1.004 & --- \\
Average CUE & 0.461 & kgCO$_2$/kWh$_{IT}$ \\
Total Carbon Emissions & 65{,}108 & kg CO$_2$ \\
Water Usage (Total) & 0 & Litres \\
Total Annual OpEx & \$20{,}523 & USD \\
Max Inlet Temperature & 38.5 & \degC \\
ASHRAE A1 Limit & 27.0 & \degC \\
\bottomrule
\end{tabular}
\end{table}

\section{Chilled Water System Reference Values}

\begin{table}[H]
\centering
\caption{Reference Values --- Chilled Water System}
\label{tab:ref_chilled}
\renewcommand{\arraystretch}{1.3}
\begin{tabular}{lrr}
\toprule
\textbf{Metric} & \textbf{Value} & \textbf{Unit} \\
\midrule
Total Facility Energy & 288{,}475 & kWh \\
Total IT Energy & 248{,}000 & kWh \\
Average PUE & 1.163 & --- \\
Water Usage & 515{,}331 & Litres \\
Total Carbon Emissions & 158{,}196 & kg CO$_2$ \\
Annual OpEx & \$42{,}553 & USD \\
CAPEX & \$630{,}000 & USD \\
Simple Payback Period & 50.0 & Years \\
ASHRAE Thermal Compliance & PASS & --- \\
\bottomrule
\end{tabular}
\end{table}

\newpage

%% ============================================================
\chapter{Executive Summary}
%% ============================================================

This chapter presents the complete simulation results for three data center cooling strategies
evaluated under identical AI-era workload conditions over a full 8{,}760-hour annual cycle.
The simulation framework integrates CloudSim Plus 8.x Discrete Event Simulation (DES) with
custom-built physics engines implemented in Java 17 / Spring Boot 3.x.

\section{Key Findings at a Glance}

\begin{table}[H]
\centering
\caption{Annual Performance Summary --- All Three Cooling Techniques}
\label{tab:exec_summary}
\renewcommand{\arraystretch}{1.4}
\begin{tabular}{lrrrr}
\toprule
\textbf{Metric} & \textbf{Air-Side Econ.} & \textbf{Evaporative} & \textbf{Chilled Water} & \textbf{Best} \\
\midrule
Total Energy (kWh/yr) & 478{,}090 & 141{,}539 & 288{,}475 & Evaporative \\
IT Energy (kWh/yr) & 372{,}125 & 141{,}173 & 248{,}000 & Evaporative \\
Average PUE & 1.285 & 1.002 & 1.163 & Evaporative \\
Average CUE (kgCO$_2$/kWh$_{IT}$) & 0.591 & 0.461 & 0.549 & Evaporative \\
Annual OpEx (USD) & \$97{,}033 & \$20{,}523 & \$42{,}553 & Evaporative \\
Carbon Emissions (kg/yr) & 219{,}921 & 65{,}108 & 158{,}196 & Evaporative \\
Water Usage (L/yr) & 0 & 0 & 515{,}331 & Air-Side / Evap \\
CAPEX (USD) & \$45{,}000 & Low & \$630{,}000 & Evaporative \\
Payback Period (yrs) & 491.9 & N/A & 50.0 & --- \\
ASHRAE Thermal Compliance & Marginal & FAIL & PASS & Chilled Water \\
\bottomrule
\end{tabular}
\end{table}

\noindent\textbf{Overall ML Recommendation:} The Random Forest recommender system recommends
\textbf{Evaporative Cooling} as the optimal strategy for the simulated conditions (model confidence: 94.75\%),
achieving the lowest PUE (1.002), lowest annual OpEx (\$20{,}523), and lowest carbon emissions
(65{,}108 kg CO$_2$/yr). However, it carries a critical thermal compliance failure --- rack inlet
temperatures reached 38.5\degC, exceeding the ASHRAE Class A1 limit of 27\degC --- making it
unsuitable as a standalone solution for high-density AI workloads without supplemental mechanical cooling.
The \textbf{Chilled Water System} is the only technique achieving full ASHRAE thermal compliance.

\newpage

%% ============================================================
\chapter{Simulation Architecture \& Co-Simulation Design}
%% ============================================================

\section{System Architecture Overview}

The simulation platform is built on a lock-step co-simulation model that tightly couples a
Discrete Event Simulation (DES) layer with a physics-based thermal engine.
Figure~\ref{fig:cosim_arch} illustrates the overall data flow.

\begin{figure}[H]
\centering
\fbox{\parbox{0.85\textwidth}{\centering\vspace{1cm}
\textit{[Figure: Co-Simulation Architecture Diagram ---\\
CloudSim Plus DES $\leftrightarrow$ Physics Engine $\leftrightarrow$ ML Recommender pipeline]}
\vspace{1cm}}}
\caption{Co-Simulation Architecture: CloudSim Plus DES coupled with three physics engines}
\label{fig:cosim_arch}
\end{figure}

\begin{table}[H]
\centering
\caption{Co-Simulation Layer Responsibilities}
\label{tab:arch_layers}
\renewcommand{\arraystretch}{1.3}
\begin{tabular}{p{3.5cm}p{5cm}p{5.5cm}}
\toprule
\textbf{Layer} & \textbf{Technology} & \textbf{Responsibility} \\
\midrule
Computational Layer & CloudSim Plus 8.x & VM provisioning, task scheduling, hourly $P_{IT}$ output \\
Thermal Layer & Spring Boot 3.x / Java 17 & Physics engine: airflow, heat balance, COP calculations \\
Synchronisation & Lock-step clock ($\Delta t = 3600$\,s) & Pulls $P_{IT}$ each tick; returns $P_{cooling}$ to global tracker \\
Economic Layer & ProjectionEngine.java & 5-year NPV, carbon tax escalation, TOU tariffs \\
ML Recommender & Random Forest (scikit-learn) & Ranks techniques by composite score \\
\bottomrule
\end{tabular}
\end{table}

\section{Workload Model}

All three simulations use an identical AI-aware composite workload:
\begin{equation}
U_{total}(t) = \bigl[U_{diurnal}(t) \times M_{AI}\bigr] + \sigma(t)
\label{eq:workload}
\end{equation}
where $U_{diurnal}$ is a sinusoidal function peaking at 14:00, $M_{AI} = 1.3$ (Mixed AI workload),
and $\sigma(t)$ is Gaussian noise ($\pm 5\%$). Weekend load is reduced by 70\%.
UPS and PDU losses are modelled with quadratic efficiency curves.

\section{Common Simulation Parameters}

\begin{table}[H]
\centering
\caption{Common Simulation Input Parameters}
\label{tab:sim_params}
\renewcommand{\arraystretch}{1.3}
\begin{tabular}{llll}
\toprule
\textbf{Category} & \textbf{Parameter} & \textbf{Value} & \textbf{Source} \\
\midrule
Thermal & ASHRAE Max Inlet ($T_{max}$) & 27.0\degC & ASHRAE TC 9.9 \\
Thermal & Thermal Mass (Rack) & 15.0 kJ/K & Design Default \\
Psychrometric & Air Density ($\rho$) & 1.2 kg/m$^3$ & Ideal Gas Law \\
Psychrometric & Latent Heat ($L_v$) & 2{,}260 kJ/kg & Water Vaporisation \\
Power & Chiller Reference COP & 3.0 & Conservative Baseline \\
Economic & Base Electricity Rate & \$0.12/kWh & Regional Default \\
Carbon & Carbon Tax (2030) & \$254/ton CO$_2$ & IPCC Pathway \\
Workload & AI Training Multiplier & 1.8$\times$ & Sustained High Load \\
Control & Mode Switch Hysteresis & 2.0\degC & Stability Buffer \\
Simulation & Total Hours & 8{,}760 & Full Annual Cycle \\
\bottomrule
\end{tabular}
\end{table}

\newpage

%% ============================================================
\chapter{Air-Side Economizer: Detailed Results}
\label{ch:air}
%% ============================================================

\section{System Description}

The Air-Side Economizer leverages ambient outdoor air to offset mechanical cooling.
A decision-tree state machine selects between three operating modes based on outdoor
dry-bulb temperature ($T_{out}$) and relative humidity ($RH_{out}$):

\begin{itemize}[noitemsep]
  \item \textbf{FULL\_ECON}: $T_{out} \leq 24$\degC \textit{and} $RH_{out} \leq 60\%$ --- 100\% outdoor air, compressor off.
  \item \textbf{PARTIAL\_TRIM}: $T_{out} \leq 24$\degC \textit{but} $RH_{out} > 60\%$ --- 20\% outdoor air, partial mechanical assist.
  \item \textbf{MECHANICAL\_ONLY}: $T_{out} > 24$\degC --- 0\% outdoor air, full compressor operation.
\end{itemize}

A 15\% filter penalty is applied to fan power in economiser modes (MERV filtration resistance).
Required volumetric airflow is:
\begin{equation}
V_{req}\ (\text{CFM}) = \frac{P_{IT,kW} \times 3160}{\rho \cdot C_p \cdot \Delta T}
\label{eq:air_cfm}
\end{equation}

\section{Annual Summary Results}

\subsection{Energy and Emissions}

\begin{table}[H]
\centering
\caption{Air-Side Economizer --- Annual Energy \& Emissions Summary}
\label{tab:air_annual}
\renewcommand{\arraystretch}{1.3}
\begin{tabular}{lrr}
\toprule
\textbf{Metric} & \textbf{Value} & \textbf{Unit} \\
\midrule
Total IT Energy & 372{,}125 & kWh \\
Total Cooling Energy & 105{,}965 & kWh \\
Total Facility Energy & 478{,}090 & kWh \\
Average PUE & 1.285 & --- \\
Average CUE & 0.591 & kgCO$_2$/kWh$_{IT}$ \\
Total Carbon Emissions & 219{,}921 & kg CO$_2$ \\
Carbon Savings vs. Baseline & 88{,}198 & kg CO$_2$ \\
Energy Savings vs. Baseline & 28.6\% & --- \\
Water Usage & 0 & Litres \\
\bottomrule
\end{tabular}
\end{table}

\noindent\textbf{Interpretation:} The air-side economizer achieved a 28.62\% energy saving compared
to a baseline mechanical-only system. The average PUE of 1.285 indicates that for every 1 kW of IT load,
the facility consumed 1.285 kW total --- the overhead being fan power and residual mechanical cooling.
The zero water usage is a key advantage over evaporative and chilled water systems.

\subsection{Economic Results}

\begin{table}[H]
\centering
\caption{Air-Side Economizer --- Economic Analysis}
\label{tab:air_econ}
\renewcommand{\arraystretch}{1.3}
\begin{tabular}{lrr}
\toprule
\textbf{Metric} & \textbf{Value} & \textbf{Unit} \\
\midrule
CAPEX & \$45{,}000 & USD \\
Annual Electricity Cost & \$69{,}323 & USD \\
Annual Carbon Tax Cost & \$27{,}710 & USD \\
Total Annual OpEx & \$97{,}033 & USD \\
Annual Savings vs. Baseline & \$91 & USD \\
Simple Payback Period & 491.9 & Years \\
\bottomrule
\end{tabular}
\end{table}

\noindent\textbf{Note on Payback Period:} The extremely long payback period (491.9 years) reflects
the high carbon tax burden (\$27{,}710/yr) applied at the IPCC 2030 rate of \$254/ton CO$_2$.
The system's carbon emissions (219{,}921 kg/yr) are high because the simulation site uses a
carbon-intensive grid (0.55 kgCO$_2$/kWh).

\section{Hourly Results --- First 24 Hours (Representative Sample)}

\begin{table}[H]
\centering
\caption{Air-Side Economizer --- Hourly Results (Hours 0--5)}
\label{tab:air_hourly}
\renewcommand{\arraystretch}{1.2}
\scriptsize
\begin{tabular}{rrrrrrrrrr}
\toprule
\textbf{Hr} & \textbf{T$_{out}$} & \textbf{RH} & \textbf{IT Load} & \textbf{CFM} & \textbf{Mode} & \textbf{Fan} & \textbf{Mech} & \textbf{Total} & \textbf{PUE} \\
 & (\degC) & (\%) & (kW) & & & (kW) & (kW) & (kW) & \\
\midrule
0 & 24.00 & 60.00 & 35.00 & 7634 & FULL\_ECON & 5.27 & 2.97 & 43.23 & 1.235 \\
1 & 24.01 & 59.99 & 34.13 & 7446 & FULL\_ECON & 5.14 & 2.91 & 42.18 & 1.236 \\
2 & 24.01 & 59.98 & 33.84 & 7382 & FULL\_ECON & 5.09 & 2.88 & 41.82 & 1.236 \\
3 & 24.02 & 59.97 & 34.13 & 7446 & FULL\_ECON & 5.14 & 2.92 & 42.19 & 1.236 \\
4 & 24.02 & 59.96 & 35.00 & 7634 & FULL\_ECON & 5.27 & 3.00 & 43.26 & 1.236 \\
5 & 24.03 & 59.95 & 36.37 & 7934 & FULL\_ECON & 5.47 & 3.13 & 44.97 & 1.237 \\
\bottomrule
\end{tabular}
\end{table}

\noindent\textbf{Key Observation:} Hours 0--5 operate cleanly in FULL\_ECON mode with no violations.
During peak workload hours (10--18), the system breaches the 10{,}000 CFM airflow limit as IT load peaks,
generating engineering recommendations to upgrade to liquid cooling (direct-to-chip) for AI-density workloads.

\section{Parameter-by-Parameter Explanation}

\begin{table}[H]
\centering
\caption{Air-Side Economizer --- Simulation Output Parameters Explained}
\label{tab:air_params}
\renewcommand{\arraystretch}{1.3}
\scriptsize
\begin{tabular}{p{3cm}p{2.5cm}p{4.5cm}p{4.5cm}}
\toprule
\textbf{Parameter} & \textbf{Appears In} & \textbf{What It Means} & \textbf{How Determined / Formula} \\
\midrule
\texttt{itLoad\_kW} & hourlyResults & CloudSim-derived IT power at each hour & CloudSim DES: $P_{IT}(t) = \sum_{i} P_{VM,i}(t)$ \\
\texttt{requiredAirflow\_CFM} & hourlyResults & Air volume needed to remove generated heat & Eq.~\ref{eq:air_cfm}: $V_{req} = \frac{P_{IT} \times 3160}{\rho C_p \Delta T}$ \\
\texttt{airflowViolation} & hourlyResults & Flag for airflow demand exceeding configured envelope & Boolean: $V_{req} > V_{max}$ (10{,}000 CFM limit) \\
\texttt{mode} & hourlyResults & Economizer control state selected for that hour & Decision tree: FULL\_ECON / PARTIAL\_TRIM / MECHANICAL\_ONLY \\
\texttt{q\_free\_kW} & hourlyResults & Heat removed by outdoor-air free cooling & $Q_{free} = \text{oaFraction} \times \rho \times (V/2118.88) \times C_p \times \Delta T_{free}$ \\
\texttt{mech\_load\_kW} & hourlyResults & Remaining heat handled by mechanical cooling & $Q_{mech} = \max(Q_{IT} - Q_{free}, 0)$ \\
\texttt{fanPower\_kW} & hourlyResults & Fan electrical power needed for airflow delivery & $P_{fan} = (V_{CFM} \times \text{fanEff} \times \text{filterFactor}) / 1000$ \\
\texttt{mechPower\_kW} & hourlyResults & Mechanical compressor power & $P_{mech} = Q_{mech} / \text{COP}_{mech}$ \\
\texttt{totalPower\_kW} & hourlyResults & Combined IT plus cooling power requirement & $P_{total} = P_{IT} + P_{fan} + P_{mech}$ \\
\texttt{pue} & hourlyResults & Facility efficiency ratio for each hour & $\text{PUE} = P_{total} / P_{IT}$ \\
\texttt{cue} & hourlyResults & Carbon intensity indicator tied to hourly operation & $\text{CUE} = (P_{total} \times \text{gridCarbonFactor}) / P_{IT}$ \\
\texttt{violationMsg} & hourlyResults & Generated engineering warning text & String: violation details + recommendation \\
\texttt{paybackPeriodYears} & summary & Economic recovery horizon of strategy & $\text{Payback} = \text{CAPEX} / \text{annualSavings}$ \\
\bottomrule
\end{tabular}
\end{table}

\section{Frontend Graphs Displayed}

The following graphs are displayed on the frontend for Air-Side Economization results:

\begin{enumerate}[noitemsep]
  \item \textbf{Hourly Power Breakdown} (Figure~\ref{fig:air_power}): Shows IT load, fan power, mechanical power, and total power over 24 hours or full year. Reveals periods where non-IT overhead expands and control strategy loses efficiency.
  
  \item \textbf{Airflow vs Free/Mechanical Cooling} (Figure~\ref{fig:air_airflow}): Plots required airflow (CFM) against free cooling and mechanical cooling contributions. Confirms when airflow ceiling drives mechanical fallback.
  
  \item \textbf{Cooling Mode Distribution} (Figure~\ref{fig:air_modes}): Pie chart or bar chart showing annual operating mode share (FULL\_ECON, PARTIAL\_TRIM, MECHANICAL\_ONLY). Explains operational dependence on favorable outdoor windows.
  
  \item \textbf{Full Cost Structure} (Figure~\ref{fig:air_cost}): Stacked bar chart of electricity cost, carbon tax cost, and total OpEx. Identifies which operational elements dominate budget pressure.
  
  \item \textbf{Rack Analysis (CloudSim)} (Figure~\ref{fig:air_racks}): Heatmap or bar chart of rack-level workload response. Shows hotspot concentration and airflow risk under AI bursts.
  
  \item \textbf{Hourly PUE and CUE} (Figure~\ref{fig:air_pue_cue}): Line charts of PUE and CUE over time. Captures variability under weather and workload fluctuations.
  
  \item \textbf{Ambient Conditions} (Figure~\ref{fig:air_ambient}): Line charts of outdoor temperature and humidity profile. Provides causal context for mode switching behavior.
  
  \item \textbf{Performance Radar} (Figure~\ref{fig:air_radar}): Radar chart with normalized multi-metric profile (efficiency, cost, emissions, thermal compliance). Summarizes balance between efficiency, cost, and sustainability.
  
  \item \textbf{5-Year Projection} (Figure~\ref{fig:air_5yr}): Line chart showing forward cost/emissions pattern. Shows long-term exposure under repeated annual behavior.
\end{enumerate}

\begin{figure}[H]
\centering
\fbox{\parbox{0.75\textwidth}{\centering\vspace{2cm}
\textit{[Figure: Hourly Power Breakdown --- Air-Side Economizer]}
\vspace{2cm}}}
\caption{Hourly Power Breakdown for Air-Side Economizer (8{,}760 hours)}
\label{fig:air_power}
\end{figure}

\begin{figure}[H]
\centering
\fbox{\parbox{0.75\textwidth}{\centering\vspace{2cm}
\textit{[Figure: Airflow vs Free/Mechanical Cooling --- Air-Side Economizer]}
\vspace{2cm}}}
\caption{Airflow Requirement vs Free/Mechanical Cooling Split}
\label{fig:air_airflow}
\end{figure}

\begin{figure}[H]
\centering
\fbox{\parbox{0.75\textwidth}{\centering\vspace{2cm}
\textit{[Figure: Cooling Mode Distribution --- Air-Side Economizer]}
\vspace{2cm}}}
\caption{Annual Cooling Mode Distribution (FULL\_ECON / PARTIAL\_TRIM / MECHANICAL\_ONLY)}
\label{fig:air_modes}
\end{figure}

\begin{figure}[H]
\centering
\fbox{\parbox{0.75\textwidth}{\centering\vspace{2cm}
\textit{[Figure: Full Cost Structure --- Air-Side Economizer]}
\vspace{2cm}}}
\caption{Annual Cost Structure (Electricity + Carbon Tax)}
\label{fig:air_cost}
\end{figure}

\begin{figure}[H]
\centering
\fbox{\parbox{0.75\textwidth}{\centering\vspace{2cm}
\textit{[Figure: Rack Analysis (CloudSim) --- Air-Side Economizer]}
\vspace{2cm}}}
\caption{Rack-Level Workload Distribution and Hotspot Analysis}
\label{fig:air_racks}
\end{figure}

\begin{figure}[H]
\centering
\fbox{\parbox{0.75\textwidth}{\centering\vspace{2cm}
\textit{[Figure: Hourly PUE and CUE --- Air-Side Economizer]}
\vspace{2cm}}}
\caption{Hourly PUE and CUE Trends Over 8{,}760 Hours}
\label{fig:air_pue_cue}
\end{figure}

\begin{figure}[H]
\centering
\fbox{\parbox{0.75\textwidth}{\centering\vspace{2cm}
\textit{[Figure: Ambient Conditions --- Air-Side Economizer]}
\vspace{2cm}}}
\caption{Outdoor Temperature and Humidity Profile (8{,}760 hours)}
\label{fig:air_ambient}
\end{figure}

\begin{figure}[H]
\centering
\fbox{\parbox{0.75\textwidth}{\centering\vspace{2cm}
\textit{[Figure: Performance Radar --- Air-Side Economizer]}
\vspace{2cm}}}
\caption{Multi-Criteria Performance Radar Chart}
\label{fig:air_radar}
\end{figure}

\begin{figure}[H]
\centering
\fbox{\parbox{0.75\textwidth}{\centering\vspace{2cm}
\textit{[Figure: 5-Year Projection --- Air-Side Economizer]}
\vspace{2cm}}}
\caption{5-Year Cost and Emissions Projection}
\label{fig:air_5yr}
\end{figure}

\section{ML Recommendation Response for Air-Side Economizer}

The ML recommender system evaluated Air-Side Economizer with the following output:

\begin{table}[H]
\centering
\caption{ML Recommendation --- Air-Side Economizer}
\label{tab:air_ml}
\renewcommand{\arraystretch}{1.3}
\begin{tabular}{lr}
\toprule
\textbf{Metric} & \textbf{Value} \\
\midrule
Model Probability & 2.08\% \\
Feasibility & True \\
Composite Score & 0.800 \\
Annual Cost & \$97{,}033 \\
Annual Emissions & 219{,}921 kg CO$_2$ \\
Annual Water Usage & 0 L \\
Violations & 0 \\
\bottomrule
\end{tabular}
\end{table}

\noindent\textbf{ML Interpretation:} The ML model assigns a low probability (2.08\%) to Air-Side Economizer
due to its high annual cost (\$97{,}033) and high carbon emissions (219{,}921 kg CO$_2$/yr) compared to
Evaporative Cooling. The composite score of 0.800 reflects moderate performance across cost, emissions,
and water usage metrics.

\newpage

%% ============================================================
\chapter{Evaporative Cooling: Detailed Results}
\label{ch:evap}
%% ============================================================

\section{System Description}

The evaporative cooling system leverages the latent heat of vaporisation ($L_v = 2{,}260$ kJ/kg)
to reduce supply air temperature. The simulation implements \textbf{Direct Evaporative Cooling (DEC)}
and \textbf{Indirect Evaporative Cooling (IEC)} modes. Supply temperature is modelled as:
\begin{equation}
T_{supply} = T_{db} - \eta_{actual} \times (T_{db} - T_{wb})
\label{eq:evap_supply}
\end{equation}
where $\eta_{actual}$ is the saturation effectiveness (60--95\%), $T_{db}$ is the outdoor dry-bulb
temperature, and $T_{wb}$ is the wet-bulb temperature. A lumped capacitance model prevents thermal spikes:
\begin{equation}
C_{total}\frac{dT}{dt} = Q_{IT}(t) - Q_{cooling}(t)
\label{eq:evap_thermal_mass}
\end{equation}
where $C_{total} = C_{rack} + C_{enclosure} = 15 + 30 = 45$ kJ/K.

\section{Annual Summary Results}

\subsection{Energy and Emissions}

\begin{table}[H]
\centering
\caption{Evaporative Cooling --- Annual Energy \& Emissions Summary}
\label{tab:evap_annual}
\renewcommand{\arraystretch}{1.3}
\begin{tabular}{lrr}
\toprule
\textbf{Metric} & \textbf{Value} & \textbf{Unit} \\
\midrule
Total IT Energy & 141{,}173 & kWh \\
Fan Energy (Auxiliary) & 366 & kWh \\
DX Compressor Energy & 0 & kWh \\
Pump Energy & 0 & kWh \\
Total Facility Energy & 141{,}539 & kWh \\
Average PUE & \textbf{1.002} & --- \\
Peak PUE & 1.004 & --- \\
Average CUE & 0.461 & kgCO$_2$/kWh$_{IT}$ \\
Total Carbon Emissions & 65{,}108 & kg CO$_2$ \\
Water Usage (Total) & 0 & Litres \\
Cooling Failure Hours & 0 & Hours \\
\bottomrule
\end{tabular}
\end{table}

\noindent The evaporative system achieves a near-ideal PUE of \textbf{1.002}, meaning only 0.2\%
overhead above the IT load --- the best result of all three techniques. Zero DX compressor energy
is consumed, and the system operates entirely on fan power.

\subsection{Water and Cost}

\begin{table}[H]
\centering
\caption{Evaporative Cooling --- Water and Economic Summary}
\label{tab:evap_water_econ}
\renewcommand{\arraystretch}{1.3}
\begin{tabular}{lrr}
\toprule
\textbf{Metric} & \textbf{Value} & \textbf{Unit} \\
\midrule
Annual Electricity Cost & \$20{,}523 & USD \\
Water Cost & \$0 & USD \\
Total Annual OpEx & \textbf{\$20{,}523} & USD \\
OpEx per kWh$_{IT}$ & \$0.1454 & USD/kWh \\
OpEx per Server (Annual) & \$213.78 & USD \\
CO$_2$ per kWh$_{IT}$ & 0.461 & kg/kWh \\
CO$_2$ per Server (Annual) & 678.21 & kg \\
\bottomrule
\end{tabular}
\end{table}

\section{Hourly Results --- First 6 Hours (Representative Sample)}

\begin{table}[H]
\centering
\caption{Evaporative Cooling --- Hourly Results (Hours 0--5)}
\label{tab:evap_hourly}
\renewcommand{\arraystretch}{1.2}
\scriptsize
\begin{tabular}{rrrrrrrrrr}
\toprule
\textbf{Hr} & \textbf{T$_{amb}$} & \textbf{RH} & \textbf{IT Load} & \textbf{Fan} & \textbf{Total} & \textbf{PUE} & \textbf{T$_{inlet}$} & \textbf{T$_{supply}$} & \textbf{Mode} \\
 & (\degC) & (\%) & (kW) & (kW) & (kW) & & (\degC) & (\degC) & \\
\midrule
0 & 25.0 & 50.0 & 11.760 & 0.0142 & 11.774 & 1.0012 & 33.87 & 18.87 & DEC \\
1 & 25.0 & 50.0 & 11.760 & 0.0142 & 11.774 & 1.0012 & 33.88 & 18.88 & DEC \\
2 & 25.0 & 50.0 & 11.760 & 0.0142 & 11.774 & 1.0012 & 33.88 & 18.88 & DEC \\
3 & 25.0 & 50.0 & 11.760 & 0.0142 & 11.774 & 1.0012 & 33.89 & 18.89 & DEC \\
4 & 25.0 & 50.0 & 11.760 & 0.0142 & 11.774 & 1.0012 & 33.87 & 18.87 & DEC \\
5 & 25.0 & 50.0 & 11.798 & 0.0144 & 11.813 & 1.0012 & 33.88 & 18.88 & DEC \\
\bottomrule
\end{tabular}
\end{table}

\noindent\textbf{Critical Observation:} Despite the excellent PUE of 1.0012, the rack inlet temperature
($T_{inlet}$) is consistently around 33.87--33.89\degC, which is \textbf{6.87\degC above the ASHRAE
Class A1 limit of 27\degC}. This is the fundamental thermal compliance failure of standalone evaporative
cooling for high-density AI workloads.

\section{Cooling Assessment and Thermal Compliance}

\begin{table}[H]
\centering
\caption{Evaporative Cooling --- Cooling Assessment Checks}
\label{tab:evap_assess}
\renewcommand{\arraystretch}{1.3}
\begin{tabular}{llp{8cm}}
\toprule
\textbf{Check} & \textbf{Result} & \textbf{Detail} \\
\midrule
Heat Balance & \fail & Cooling capacity 3.5\% below IT heat load at peak \\
Inlet Temperature & \fail & Max inlet temp: 38.5\degC (ASHRAE limit: 27\degC) \\
Humidity & \pass & Supply humidity within acceptable range \\
Energy Efficiency & \pass & PUE 1.002 --- excellent \\
\midrule
\textbf{Overall Status} & \multicolumn{2}{l}{\textcolor{failred}{\textbf{INSUFFICIENT\_COOLING}} (confidence: 95\%)} \\
\bottomrule
\end{tabular}
\end{table}

\noindent\textbf{Critical Finding:} Rack inlet temperatures reached \textbf{38.5\degC}, exceeding
the ASHRAE Class A1 limit of 27\degC by 11.5\degC. This triggers the thermal throttling model:
\begin{equation}
\Phi_{adj} = \Phi_{base} \times (1 - \gamma \cdot \Delta T)
\label{eq:throttle}
\end{equation}
At $\Delta T = 11.5$\degC and $\gamma = 0.10$ (aggressive throttling above 29\degC), computational
throughput is reduced by up to \textbf{115\%} --- effectively a full shutdown.

\section{Parameter-by-Parameter Explanation}

\begin{table}[H]
\centering
\caption{Evaporative Cooling --- Simulation Output Parameters Explained}
\label{tab:evap_params}
\renewcommand{\arraystretch}{1.3}
\scriptsize
\begin{tabular}{p{3cm}p{2.5cm}p{4.5cm}p{4.5cm}}
\toprule
\textbf{Parameter} & \textbf{Appears In} & \textbf{What It Means} & \textbf{How Determined / Formula} \\
\midrule
\texttt{itLoadKW} & hourly\_data & CloudSim-generated compute heat load per hour & Lock-step DES: $P_{IT}(t) = \text{CloudSim.runFor}(3600)$ \\
\texttt{coolingMode} & hourly\_data & Evaporative operating mode selected by controller & DEC or IEC based on ambient conditions \\
\texttt{ambientTempC} & hourly\_data & Outdoor dry-bulb condition used in simulation & Live weather API (hourly TMY3 data) \\
\texttt{ambientHumidity} & hourly\_data & Outdoor humidity condition for psychrometrics & Live weather API; drives wet-bulb depression \\
\texttt{supplyTempC} & hourly\_data & Delivered supply-air temperature & Eq.~\ref{eq:evap_supply}: $T_{supply} = T_{db} - \eta(T_{db}-T_{wb})$ \\
\texttt{supplyHumidity} & hourly\_data & Delivered supply-air humidity & Psychrometric calculation from supply conditions \\
\texttt{inletTempC} & hourly\_data & Estimated rack inlet thermal condition & $T_{inlet} = T_{supply} + \Delta T_{bypass}$ (thermal mass model) \\
\texttt{coolingCapacityKW} & hourly\_data & Cooling produced by evaporative path & $Q_{cooling} = \dot{V} \times \rho_{air} \times C_p \times \Delta T$ \\
\texttt{fanPowerKW} & hourly\_data & Fan power needed for evaporative airflow & Affinity law: $P_{fan} \propto \text{speed}^3$ \\
\texttt{dxPowerKW} & hourly\_data & DX compressor power (zero in pure evap mode) & $P_{DX} = 0$ when $T_{supply} \leq T_{setpoint}$ \\
\texttt{totalElectricalKW} & hourly\_data & Aggregate electrical requirement for hour & $P_{total} = P_{IT} + P_{fan} + P_{DX} + P_{pump}$ \\
\texttt{pue} & hourly\_data & Facility efficiency ratio & $\text{PUE} = P_{total} / P_{IT}$ \\
\texttt{waterEvaporationLph} & hourly\_data & Water evaporation rate (L/hr) & $\dot{m}_{evap} = Q_{cooling} \times 3600 / L_v$ \\
\texttt{cooling\_failure\_hours} & results.performance & Count of inadequate cooling events & Hours where $Q_{cooling} < Q_{IT}$ \\
\texttt{co2\_kg\_per\_kwh\_it} & results.emissions & Carbon intensity per IT energy served & $\text{CUE} = (P_{total} \times \text{gridFactor}) / P_{IT}$ \\
\texttt{pue\_average} & results.performance & Annual average PUE & $\overline{\text{PUE}} = \frac{1}{8760}\sum_{t=1}^{8760}\text{PUE}(t)$ \\
\texttt{max\_inlet\_temp\_c} & cooling\_assessment & Peak rack inlet temperature observed & $\max_{t}(T_{inlet}(t))$ over 8{,}760 hours \\
\texttt{cooling\_capacity\_avg\_kw} & cooling\_assessment & Average cooling capacity & $\overline{Q}_{cooling} = 17.11$ kW \\
\texttt{heat\_load\_avg\_kw} & cooling\_assessment & Average IT heat load & $\overline{Q}_{IT} = 17.73$ kW \\
\bottomrule
\end{tabular}
\end{table}

\section{Frontend Graphs Displayed}

The following graphs are displayed on the frontend for Evaporative Cooling results:

\begin{enumerate}[noitemsep]
  \item \textbf{Annual Evaporative Summary} (Figure~\ref{fig:evap_summary}): Top-level annual result set showing energy, cost, and emissions. Shows strong efficiency profile but must be validated against thermal adequacy.
  
  \item \textbf{Hourly Power Breakdown} (Figure~\ref{fig:evap_power}): IT and auxiliary power traces over time. Confirms low auxiliary overhead pattern during normal operation.
  
  \item \textbf{Cooling Capacity vs IT Load} (Figure~\ref{fig:evap_capacity}): Capacity-demand interaction chart. Highlights hours where thermal margin narrows under peak demand.
  
  \item \textbf{PUE vs Annual Averages} (Figure~\ref{fig:evap_pue}): Hourly efficiency vs annual baseline. Shows whether transient behavior remains near target operating range.
  
  \item \textbf{Temperature and Humidity} (Figure~\ref{fig:evap_temp_hum}): Ambient and process condition trends. Reveals environmental periods that constrain evaporative headroom.
  
  \item \textbf{Supply Air Conditions} (Figure~\ref{fig:evap_supply}): Supply/inlet temperature and humidity behavior. Checks if delivered conditions remain suitable for rack safety.
  
  \item \textbf{Cooling Mode Distribution} (Figure~\ref{fig:evap_modes}): Annual mode occupancy (DEC/IEC). Identifies operating dependence on weather profile.
  
  \item \textbf{PUE per Hour} (Figure~\ref{fig:evap_pue_hourly}): Hourly efficiency trace over 8{,}760 hours. Shows consistency and detects outlier operation windows.
  
  \item \textbf{Cooling Assessment} (Figure~\ref{fig:evap_assess_chart}): Adequacy checklist outcomes (pass/fail). Separates efficiency success from thermal compliance success.
  
  \item \textbf{Performance Radar} (Figure~\ref{fig:evap_radar}): Multi-metric normalized view. Summarizes trade-off profile quickly for decision review.
  
  \item \textbf{5-Year Projection} (Figure~\ref{fig:evap_5yr}): Long-horizon trend view. Shows persistence of operational advantages and risk exposure.
\end{enumerate}

\begin{figure}[H]
\centering
\fbox{\parbox{0.75\textwidth}{\centering\vspace{2cm}
\textit{[Figure: Annual Evaporative Summary]}
\vspace{2cm}}}
\caption{Annual Evaporative Cooling Summary Dashboard}
\label{fig:evap_summary}
\end{figure}

\begin{figure}[H]
\centering
\fbox{\parbox{0.75\textwidth}{\centering\vspace{2cm}
\textit{[Figure: Hourly Power Breakdown --- Evaporative Cooling]}
\vspace{2cm}}}
\caption{Hourly Power Breakdown for Evaporative Cooling}
\label{fig:evap_power}
\end{figure}

\begin{figure}[H]
\centering
\fbox{\parbox{0.75\textwidth}{\centering\vspace{2cm}
\textit{[Figure: Cooling Capacity vs IT Load --- Evaporative Cooling]}
\vspace{2cm}}}
\caption{Cooling Capacity vs IT Load (Thermal Margin Analysis)}
\label{fig:evap_capacity}
\end{figure}

\begin{figure}[H]
\centering
\fbox{\parbox{0.75\textwidth}{\centering\vspace{2cm}
\textit{[Figure: PUE vs Annual Averages --- Evaporative Cooling]}
\vspace{2cm}}}
\caption{Hourly PUE vs Annual Average Baseline}
\label{fig:evap_pue}
\end{figure}

\begin{figure}[H]
\centering
\fbox{\parbox{0.75\textwidth}{\centering\vspace{2cm}
\textit{[Figure: Temperature and Humidity --- Evaporative Cooling]}
\vspace{2cm}}}
\caption{Ambient Temperature and Humidity Profile}
\label{fig:evap_temp_hum}
\end{figure}

\begin{figure}[H]
\centering
\fbox{\parbox{0.75\textwidth}{\centering\vspace{2cm}
\textit{[Figure: Supply Air Conditions --- Evaporative Cooling]}
\vspace{2cm}}}
\caption{Supply Air Temperature and Humidity at Rack Inlet}
\label{fig:evap_supply}
\end{figure}

\begin{figure}[H]
\centering
\fbox{\parbox{0.75\textwidth}{\centering\vspace{2cm}
\textit{[Figure: Cooling Mode Distribution --- Evaporative Cooling]}
\vspace{2cm}}}
\caption{Annual Cooling Mode Distribution (DEC / IEC)}
\label{fig:evap_modes}
\end{figure}

\begin{figure}[H]
\centering
\fbox{\parbox{0.75\textwidth}{\centering\vspace{2cm}
\textit{[Figure: PUE per Hour --- Evaporative Cooling]}
\vspace{2cm}}}
\caption{Hourly PUE Trace Over 8{,}760 Hours}
\label{fig:evap_pue_hourly}
\end{figure}

\begin{figure}[H]
\centering
\fbox{\parbox{0.75\textwidth}{\centering\vspace{2cm}
\textit{[Figure: Cooling Assessment --- Evaporative Cooling]}
\vspace{2cm}}}
\caption{Cooling Assessment Pass/Fail Checklist}
\label{fig:evap_assess_chart}
\end{figure}

\begin{figure}[H]
\centering
\fbox{\parbox{0.75\textwidth}{\centering\vspace{2cm}
\textit{[Figure: Performance Radar --- Evaporative Cooling]}
\vspace{2cm}}}
\caption{Multi-Criteria Performance Radar Chart}
\label{fig:evap_radar}
\end{figure}

\begin{figure}[H]
\centering
\fbox{\parbox{0.75\textwidth}{\centering\vspace{2cm}
\textit{[Figure: 5-Year Projection --- Evaporative Cooling]}
\vspace{2cm}}}
\caption{5-Year Cost and Emissions Projection}
\label{fig:evap_5yr}
\end{figure}

\section{ML Recommendation Response for Evaporative Cooling}

\begin{table}[H]
\centering
\caption{ML Recommendation --- Evaporative Cooling}
\label{tab:evap_ml}
\renewcommand{\arraystretch}{1.3}
\begin{tabular}{lr}
\toprule
\textbf{Metric} & \textbf{Value} \\
\midrule
Model Probability & \textbf{94.75\%} \\
Feasibility & True \\
Composite Score & \textbf{0.000} (lowest = best) \\
Annual Cost & \$20{,}523 \\
Annual Emissions & 65{,}108 kg CO$_2$ \\
Annual Water Usage & 0 L \\
Violations & 0 \\
\bottomrule
\end{tabular}
\end{table}

\noindent\textbf{ML Interpretation:} The Random Forest model assigns the highest probability (94.75\%)
to Evaporative Cooling because it achieves the lowest composite score (0.000) across cost, emissions,
and water usage metrics. The model recommends Evaporative as the primary technique, noting that it
saves \$22{,}030/yr vs Chilled Water and \$76{,}510/yr vs Air-Side Economizer. The model's
\texttt{future\_impact\_paragraph} projects: Year 1 cost \$20{,}523, Year 3 cumulative \$61{,}569,
Year 5 cumulative \$102{,}616.

\noindent\textbf{Important Caveat:} The ML model does not directly penalise thermal compliance failures
in its scoring function. The thermal compliance failure (inlet temp 38.5\degC) is captured in the
\texttt{cooling\_assessment} block of the simulation response but is not a hard constraint in the
Random Forest feature set. This is a known limitation of the current ML pipeline.

\newpage

%% ============================================================
\chapter{Chilled Water System: Detailed Results}
\label{ch:chw}
%% ============================================================

\section{System Description}

The chilled water system uses the Energy Input Ratio (EIR) framework for plant operations.
Efficiency is a function of chilled water supply temperature ($T_{chw}$) and condenser water
temperature ($T_{cond}$). Variable-speed pumping and cooling tower fans follow affinity laws
where power scales cubically with speed ($P \propto \text{speed}^3$). The system includes:

\begin{itemize}[noitemsep]
  \item Centrifugal chillers with EIR-based part-load performance curves
  \item Variable primary/secondary pumping loops
  \item Cooling tower with wet-bulb approach control
  \item CRAH units for in-row heat rejection
  \item Bayesian calibration for real-world performance matching
\end{itemize}

The chiller COP is computed via the EIR method:
\begin{equation}
\text{COP}_{actual} = \frac{\text{COP}_{ref}}{\text{EIR}_{FPLR} \times \text{EIR}_{FT}}
\label{eq:chw_cop}
\end{equation}
where $\text{EIR}_{FPLR}$ is the part-load ratio correction and $\text{EIR}_{FT}$ is the
temperature correction factor. Chiller power is:
\begin{equation}
P_{chiller} = \frac{Q_{cooling}}{\text{COP}_{actual}}
\label{eq:chw_power}
\end{equation}

\section{Annual Summary Results}

\subsection{Energy and Emissions}

\begin{table}[H]
\centering
\caption{Chilled Water System --- Annual Energy \& Emissions Summary}
\label{tab:chw_annual}
\renewcommand{\arraystretch}{1.3}
\begin{tabular}{lrr}
\toprule
\textbf{Metric} & \textbf{Value} & \textbf{Unit} \\
\midrule
Total Facility Energy & 288{,}475 & kWh \\
Cooling Load & 40{,}487 & kWh \\
Average PUE & 1.163 & --- \\
Water Usage & 515{,}331 & Litres \\
WUE & 1.786 & L/kWh$_{IT}$ \\
Average COP & 1.000 & --- \\
Peak Cooling Load & 5.853 & kW \\
Total Carbon Emissions & 158{,}196 & kg CO$_2$ \\
Annual OpEx & \$42{,}553 & USD \\
\bottomrule
\end{tabular}
\end{table}

\noindent\textbf{Note on COP:} The average COP of 1.0 reflects the simulation's conservative EIR
baseline (COP$_{ref}$ = 3.0 at design conditions). The reported value represents the effective
system COP including all auxiliary loads (pumps, cooling tower fans, CRAH units).
Real-world chillers achieve COP 4--6 at part load.

\subsection{Economic Results}

\begin{table}[H]
\centering
\caption{Chilled Water System --- Economic Analysis}
\label{tab:chw_econ}
\renewcommand{\arraystretch}{1.3}
\begin{tabular}{lrr}
\toprule
\textbf{Metric} & \textbf{Value} & \textbf{Unit} \\
\midrule
CAPEX & \$630{,}000 & USD \\
Annual OpEx & \$42{,}553 & USD \\
LCCP (15-year) & \$1{,}838{,}582 & USD \\
NPV & $-$\$1{,}092{,}355 & USD \\
Simple Payback Period & 50.0 & Years \\
\bottomrule
\end{tabular}
\end{table}

\noindent The negative NPV of $-$\$1{,}092{,}355 and 50-year payback reflect the high CAPEX
(\$630{,}000) of a full chilled water plant. This is typical for enterprise-scale deployments
where the investment is amortised over large IT loads (MW-scale).

\subsection{Phase Gate Assessment}

\begin{table}[H]
\centering
\caption{Chilled Water System --- Phase Gate Compliance}
\label{tab:chw_gates}
\renewcommand{\arraystretch}{1.3}
\begin{tabular}{llp{8cm}}
\toprule
\textbf{Gate} & \textbf{Status} & \textbf{Notes} \\
\midrule
Thermal Compliance & \pass & Inlet temps maintained within ASHRAE A2 range \\
Water Constraint & \pass & 515{,}331 L/yr within site water budget \\
Carbon Liability & \fail & 158{,}196 kg CO$_2$/yr exceeds carbon target \\
Economic Viability & \fail & NPV $-$\$1.09M; payback 50 years \\
\bottomrule
\end{tabular}
\end{table}

\noindent\textbf{Key Advantage:} The chilled water system is the \textbf{only technique to achieve
full thermal compliance} at ASHRAE Class A2 limits. It is the recommended choice for high-density
AI workloads where server inlet temperature must be guaranteed below 27\degC.

\section{Hourly Results --- First 10 Hours (Representative Sample)}

\begin{table}[H]
\centering
\caption{Chilled Water System --- Hourly Results (Hours 0--9)}
\label{tab:chw_hourly}
\renewcommand{\arraystretch}{1.2}
\scriptsize
\begin{tabular}{rrrrrrrrr}
\toprule
\textbf{Hr} & \textbf{T$_{amb}$} & \textbf{IT Load} & \textbf{Cool Load} & \textbf{Chiller} & \textbf{COP} & \textbf{Water} & \textbf{Cost} & \textbf{CO$_2$} \\
 & (\degC) & (kW) & (kW) & (kW) & & (L) & (USD) & (kg) \\
\midrule
0 & 24.00 & 21.43 & 2.757 & 2.678 & 8.000 & 43.39 & 2.636 & 14.12 \\
1 & 24.01 & 21.43 & 2.757 & 2.678 & 7.999 & 43.39 & 2.636 & 14.12 \\
2 & 24.01 & 21.43 & 2.757 & 2.679 & 7.999 & 43.39 & 2.636 & 14.12 \\
3 & 24.02 & 21.43 & 2.758 & 2.679 & 7.998 & 43.39 & 2.636 & 14.12 \\
4 & 24.02 & 21.43 & 2.758 & 2.679 & 7.998 & 43.39 & 2.636 & 14.12 \\
5 & 24.03 & 24.62 & 3.198 & 3.079 & 7.997 & 49.86 & 2.942 & 15.76 \\
6 & 24.03 & 25.09 & 3.265 & 3.138 & 7.996 & 50.81 & 4.267 & 16.00 \\
7 & 24.04 & 25.81 & 3.365 & 3.228 & 7.996 & 52.26 & 4.364 & 16.37 \\
8 & 24.05 & 26.74 & 3.498 & 3.345 & 7.995 & 54.16 & 4.493 & 16.85 \\
9 & 24.05 & 27.79 & 3.649 & 3.477 & 7.994 & 56.29 & 4.637 & 17.39 \\
\bottomrule
\end{tabular}
\end{table}

\noindent\textbf{Observation:} The chilled water system maintains stable hourly COP around 8.0 at
low ambient temperatures (24\degC), demonstrating the efficiency advantage of vapour-compression
refrigeration at near-design conditions. Water consumption scales linearly with cooling load.

\section{Parameter-by-Parameter Explanation}

\begin{table}[H]
\centering
\caption{Chilled Water System --- Simulation Output Parameters Explained}
\label{tab:chw_params}
\renewcommand{\arraystretch}{1.3}
\scriptsize
\begin{tabular}{p{3cm}p{2.5cm}p{4.5cm}p{4.5cm}}
\toprule
\textbf{Parameter} & \textbf{Appears In} & \textbf{What It Means} & \textbf{How Determined / Formula} \\
\midrule
\texttt{itLoad\_kW} & hourlyResults & CloudSim workload translated into thermal load & Lock-step DES: $P_{IT}(t) = \text{CloudSim.runFor}(3600)$ \\
\texttt{ambientTemp\_C} & hourlyResults & Ambient condition for tower and condenser response & Live weather API; influences condenser rejection \\
\texttt{coolingLoad\_kW} & hourlyResults & Cooling energy that must be supplied to match load & $Q_{cooling} = P_{IT} \times \text{coolingFraction}$ \\
\texttt{chillerPower\_kW} & hourlyResults & Compressor-side electrical consumption & Eq.~\ref{eq:chw_power}: $P_{chiller} = Q_{cooling} / \text{COP}_{actual}$ \\
\texttt{cop} & hourlyResults & Instantaneous chiller efficiency indicator & Eq.~\ref{eq:chw_cop}: EIR-based COP calculation \\
\texttt{waterUsage\_L} & hourlyResults & Hourly tower-related water consumption & $\dot{V}_{water} = Q_{rejected} \times \text{evapFraction} / \rho_w$ \\
\texttt{cost\_USD} & hourlyResults & Hourly operating cost estimate & $C_{hr} = P_{total} \times \text{tariff}(t)$ (TOU pricing) \\
\texttt{carbonEmissions\_kg} & hourlyResults & Hourly carbon impact derived from power use & $E_{CO_2} = P_{total} \times \text{gridFactor} \times 3600 / 10^6$ \\
\texttt{energyConsumption\_kWh} & results.annual & Total annual facility energy & $\sum_{t=1}^{8760} P_{total}(t)$ \\
\texttt{coolingLoad\_kWh} & results.annual & Total annual cooling energy delivered & $\sum_{t=1}^{8760} Q_{cooling}(t)$ \\
\texttt{waterUsage\_L} & results.annual & Total annual water consumption & $\sum_{t=1}^{8760} \dot{V}_{water}(t)$ \\
\texttt{pue} & results.metrics & Annual average PUE & $\overline{\text{PUE}} = E_{facility} / E_{IT}$ \\
\texttt{wue} & results.metrics & Water usage effectiveness & $\text{WUE} = V_{water,annual} / E_{IT,annual}$ \\
\texttt{averageCOP} & results.metrics & Year-level average efficiency indicator & $\overline{\text{COP}} = Q_{cooling,total} / P_{chiller,total}$ \\
\texttt{peakCoolingLoad\_kW} & results.metrics & Maximum instantaneous cooling demand & $\max_{t}(Q_{cooling}(t)) = 5.853$ kW \\
\texttt{capex\_USD} & results.economics & Capital expenditure for chilled water plant & Engineering estimate: \$630{,}000 \\
\texttt{opex\_annual\_USD} & results.economics & Annual operating expenditure & $\text{OpEx} = C_{electricity} + C_{water} + C_{maintenance}$ \\
\texttt{lccp\_USD} & results.economics & Life-cycle cost of plant (15-year) & $\text{LCCP} = \text{CAPEX} + \sum_{y=1}^{15} \text{OpEx}_y / (1+r)^y$ \\
\texttt{npv\_USD} & results.economics & Net present value of investment & $\text{NPV} = -\text{CAPEX} + \sum_{y=1}^{n} \frac{\text{Savings}_y}{(1+r)^y}$ \\
\texttt{paybackPeriod\_years} & results.economics & Economic recovery horizon & $\text{Payback} = \text{CAPEX} / \text{annualSavings}$ \\
\texttt{thermalCompliance} & results.phase4Gates & ASHRAE thermal gate outcome & Rule: $T_{inlet} \leq 27$\degC for all hours \\
\texttt{waterConstraint} & results.phase4Gates & Water budget gate outcome & Rule: $V_{water} \leq V_{budget,site}$ \\
\texttt{carbonLiability} & results.phase4Gates & Carbon emissions gate outcome & Rule: $E_{CO_2} \leq E_{target}$ \\
\texttt{economicViability} & results.phase4Gates & Economic gate outcome & Rule: NPV $> 0$ or payback $< 10$ years \\
\bottomrule
\end{tabular}
\end{table}

\section{Frontend Graphs Displayed}

The following graphs are displayed on the frontend for Chilled Water System results:

\begin{enumerate}[noitemsep]
  \item \textbf{Annual Consumption Overview} (Figure~\ref{fig:chw_overview}): Annual plant summary showing energy, water, cost, and emissions. Shows balanced operating envelope with controlled thermal delivery.
  
  \item \textbf{COP Over Time} (Figure~\ref{fig:chw_cop}): Efficiency progression over 8{,}760 hours. Reveals stability of chiller efficiency under changing demand.
  
  \item \textbf{Cooling Load vs Chiller Power} (Figure~\ref{fig:chw_load_power}): Demand and plant response correlation. Confirms how consistently plant output tracks load.
  
  \item \textbf{IT Load vs Chiller Power} (Figure~\ref{fig:chw_it_chiller}): Compute-cooling coupling chart. Shows plant responsiveness to workload dynamics.
  
  \item \textbf{Water and Carbon Trend} (Figure~\ref{fig:chw_water_carbon}): Resource and emissions pattern over time. Highlights sustainability trade-off of stronger thermal control.
  
  \item \textbf{Cost Structure} (Figure~\ref{fig:chw_cost}): Annual and lifecycle cost composition. Clarifies where robustness increases cost commitment.
  
  \item \textbf{Phase 4 Compliance Gates} (Figure~\ref{fig:chw_gates}): Gate status outcomes (PASS/FAIL). Summarizes engineering readiness for production-grade thermal safety.
  
  \item \textbf{Performance Radar} (Figure~\ref{fig:chw_radar}): Normalized overall profile. Shows strong compliance profile with moderate efficiency trade-off.
  
  \item \textbf{5-Year Projection} (Figure~\ref{fig:chw_5yr}): Long-horizon behavior. Shows predictable performance trajectory under repeated demand cycles.
\end{enumerate}

\begin{figure}[H]
\centering
\fbox{\parbox{0.75\textwidth}{\centering\vspace{2cm}
\textit{[Figure: Annual Consumption Overview --- Chilled Water System]}
\vspace{2cm}}}
\caption{Annual Consumption Overview for Chilled Water System}
\label{fig:chw_overview}
\end{figure}

\begin{figure}[H]
\centering
\fbox{\parbox{0.75\textwidth}{\centering\vspace{2cm}
\textit{[Figure: COP Over Time --- Chilled Water System]}
\vspace{2cm}}}
\caption{Chiller COP Over 8{,}760 Hours}
\label{fig:chw_cop}
\end{figure}

\begin{figure}[H]
\centering
\fbox{\parbox{0.75\textwidth}{\centering\vspace{2cm}
\textit{[Figure: Cooling Load vs Chiller Power --- Chilled Water System]}
\vspace{2cm}}}
\caption{Cooling Load vs Chiller Power (Plant Response Tracking)}
\label{fig:chw_load_power}
\end{figure}

\begin{figure}[H]
\centering
\fbox{\parbox{0.75\textwidth}{\centering\vspace{2cm}
\textit{[Figure: IT Load vs Chiller Power --- Chilled Water System]}
\vspace{2cm}}}
\caption{IT Load vs Chiller Power (Compute-Cooling Coupling)}
\label{fig:chw_it_chiller}
\end{figure}

\begin{figure}[H]
\centering
\fbox{\parbox{0.75\textwidth}{\centering\vspace{2cm}
\textit{[Figure: Water and Carbon Trend --- Chilled Water System]}
\vspace{2cm}}}
\caption{Water Consumption and Carbon Emissions Trend}
\label{fig:chw_water_carbon}
\end{figure}

\begin{figure}[H]
\centering
\fbox{\parbox{0.75\textwidth}{\centering\vspace{2cm}
\textit{[Figure: Cost Structure --- Chilled Water System]}
\vspace{2cm}}}
\caption{Annual and Lifecycle Cost Structure}
\label{fig:chw_cost}
\end{figure}

\begin{figure}[H]
\centering
\fbox{\parbox{0.75\textwidth}{\centering\vspace{2cm}
\textit{[Figure: Phase 4 Compliance Gates --- Chilled Water System]}
\vspace{2cm}}}
\caption{Phase 4 Engineering Compliance Gate Results}
\label{fig:chw_gates}
\end{figure}

\begin{figure}[H]
\centering
\fbox{\parbox{0.75\textwidth}{\centering\vspace{2cm}
\textit{[Figure: Performance Radar --- Chilled Water System]}
\vspace{2cm}}}
\caption{Multi-Criteria Performance Radar Chart}
\label{fig:chw_radar}
\end{figure}

\begin{figure}[H]
\centering
\fbox{\parbox{0.75\textwidth}{\centering\vspace{2cm}
\textit{[Figure: 5-Year Projection --- Chilled Water System]}
\vspace{2cm}}}
\caption{5-Year Cost and Emissions Projection}
\label{fig:chw_5yr}
\end{figure}

\section{ML Recommendation Response for Chilled Water System}

\begin{table}[H]
\centering
\caption{ML Recommendation --- Chilled Water System}
\label{tab:chw_ml}
\renewcommand{\arraystretch}{1.3}
\begin{tabular}{lr}
\toprule
\textbf{Metric} & \textbf{Value} \\
\midrule
Model Probability & 3.17\% \\
Feasibility & True \\
Composite Score & 0.524 \\
Annual Cost & \$42{,}553 \\
Annual Emissions & 158{,}196 kg CO$_2$ \\
Annual Water Usage & 515{,}331 L \\
Violations & 0 \\
\bottomrule
\end{tabular}
\end{table}

\noindent\textbf{ML Interpretation:} The ML model assigns a low probability (3.17\%) to Chilled Water
because its cost (\$42{,}553/yr) and emissions (158{,}196 kg CO$_2$/yr) are higher than Evaporative
Cooling. However, the model notes that choosing Evaporative over Chilled Water saves approximately
\$22{,}030/yr and reduces emissions by 93{,}088 kg CO$_2$/yr while saving 515{,}331 litres of water.
The composite score of 0.524 reflects moderate performance, penalised primarily by high water usage.

\newpage

%% ============================================================
\chapter{ML Recommendation Engine: Full Response Analysis}
\label{ch:ml}
%% ============================================================

\section{ML System Overview}

The ML recommendation engine is a \textbf{Random Forest classifier} trained on 8{,}760-hour
simulation outputs from all three cooling techniques. The model takes site conditions
(ambient temperature, humidity, IT load, water stress, electricity tariff, carbon intensity)
as input features and outputs a ranked recommendation with confidence scores.

\begin{table}[H]
\centering
\caption{ML Recommender System Configuration}
\label{tab:ml_config}
\renewcommand{\arraystretch}{1.3}
\begin{tabular}{ll}
\toprule
\textbf{Parameter} & \textbf{Value} \\
\midrule
Algorithm & Random Forest Classifier \\
Library & scikit-learn (Python) \\
Training Data & Synthetic + simulation-derived dataset \\
Output Classes & AirEconomizer, Evaporative, ChilledWater \\
Scoring Metric & Composite weighted score (cost + emissions + water) \\
Feasibility Check & Constraint violation count per technique \\
\bottomrule
\end{tabular}
\end{table}

\section{Full ML Response (Authentic 8{,}760-Hour Run)}

The following is the complete ML recommendation response from \texttt{ml\_authentic\_8760\_output.json}:

\begin{table}[H]
\centering
\caption{ML Recommendation --- Full Response Summary}
\label{tab:ml_full}
\renewcommand{\arraystretch}{1.3}
\begin{tabular}{ll}
\toprule
\textbf{Field} & \textbf{Value} \\
\midrule
\texttt{current\_technique} & ChilledWater \\
\texttt{model\_recommendation} & \textbf{Evaporative} \\
\texttt{model\_confidence} & 94.75\% \\
\texttt{generated\_at\_utc} & 2026-04-07T17:33:13Z \\
\bottomrule
\end{tabular}
\end{table}

\subsection{Comparison Table from ML Response}

\begin{table}[H]
\centering
\caption{ML Recommender --- Technique Comparison Table (from JSON response)}
\label{tab:ml_comparison}
\renewcommand{\arraystretch}{1.3}
\begin{tabular}{lrrrrrr}
\toprule
\textbf{Technique} & \textbf{Feasible} & \textbf{Score} & \textbf{Annual Cost} & \textbf{Emissions (kg)} & \textbf{Water (L)} & \textbf{Violations} \\
\midrule
Evaporative & Yes & \textbf{0.000} & \$20{,}523 & 65{,}108 & 0 & 0 \\
ChilledWater & Yes & 0.524 & \$42{,}553 & 158{,}196 & 515{,}331 & 0 \\
AirEconomizer & Yes & 0.800 & \$97{,}033 & 219{,}921 & 0 & 0 \\
\bottomrule
\end{tabular}
\end{table}

\subsection{Model Probabilities}

\begin{table}[H]
\centering
\caption{Random Forest Class Probabilities}
\label{tab:ml_probs}
\renewcommand{\arraystretch}{1.3}
\begin{tabular}{lr}
\toprule
\textbf{Technique} & \textbf{Probability} \\
\midrule
Evaporative & \textbf{94.75\%} \\
ChilledWater & 3.17\% \\
AirEconomizer & 2.08\% \\
\bottomrule
\end{tabular}
\end{table}

\subsection{Justification Statements from ML Response}

The ML model generated the following justification statements:

\begin{enumerate}[noitemsep]
  \item ``Recommended technique is Evaporative, and it is feasible for this case, so it remains the final recommendation.''
  \item ``Your current technique is ChilledWater, but the suggested technique is Evaporative for better overall performance.''
  \item ``Based on your current conditions, the suggested cooling technique is Evaporative.''
  \item ``Feasibility check result: Evaporative is valid for this scenario with 0 constraint violations.''
  \item ``Compared with ChilledWater, choosing Evaporative can save about \$22{,}030 per year, reduce emissions by about 93{,}088 kg CO$_2$ per year, and save about 515{,}331 litres of water per year.''
  \item ``Compared with AirEconomizer, choosing Evaporative can save about \$76{,}510 per year, reduce emissions by about 154{,}813 kg CO$_2$ per year, and use about 0 more litres of water per year.''
\end{enumerate}

\subsection{Future Impact Projection from ML Response}

\begin{table}[H]
\centering
\caption{ML Future Impact Projection (Evaporative Cooling)}
\label{tab:ml_projection}
\renewcommand{\arraystretch}{1.3}
\begin{tabular}{lrrr}
\toprule
\textbf{Year} & \textbf{Cumulative Cost (USD)} & \textbf{Cumulative CO$_2$ (kg)} & \textbf{Cumulative Water (L)} \\
\midrule
Year 1 & \$20{,}523 & 65{,}108 & 0 \\
Year 3 & \$61{,}569 & 195{,}324 & 0 \\
Year 5 & \$102{,}616 & 325{,}540 & 0 \\
\bottomrule
\end{tabular}
\end{table}

\section{ML Parameter Explanation}

\begin{table}[H]
\centering
\caption{ML Response Parameters Explained}
\label{tab:ml_params}
\renewcommand{\arraystretch}{1.3}
\scriptsize
\begin{tabular}{p{3.5cm}p{4.5cm}p{6.5cm}}
\toprule
\textbf{Parameter} & \textbf{What It Means} & \textbf{How Determined} \\
\midrule
\texttt{model\_recommendation} & The technique the RF model predicts as optimal & $\arg\max_c P(c | \mathbf{x})$ where $\mathbf{x}$ is the feature vector \\
\texttt{model\_confidence} & Probability assigned to the recommended class & RF: fraction of decision trees voting for the winning class \\
\texttt{model\_probabilities} & Per-class probability distribution & Soft voting across all trees in the Random Forest \\
\texttt{score} & Composite weighted score (lower = better) & $s = w_1 \cdot \hat{C} + w_2 \cdot \hat{E} + w_3 \cdot \hat{W}$ (normalised cost, emissions, water) \\
\texttt{feasible} & Whether the technique passes all hard constraints & Boolean: violations $= 0$ \\
\texttt{violations} & Count of constraint violations & Checks: thermal, water, carbon, economic gates \\
\texttt{annual\_cost} & Projected annual operating cost & Directly from simulation output \\
\texttt{annual\_emissions\_kg} & Projected annual CO$_2$ emissions & Directly from simulation output \\
\texttt{annual\_water\_liters} & Projected annual water consumption & Directly from simulation output \\
\texttt{future\_impact\_paragraph} & Natural language projection summary & Template-filled from Year 1/3/5 projections \\
\texttt{why\_this\_is\_recommended} & List of justification strings & Rule-based post-processing of comparison table \\
\bottomrule
\end{tabular}
\end{table}

\newpage

%% ============================================================
\chapter{Comparative Analysis}
\label{ch:compare}
%% ============================================================

\section{Energy Efficiency Comparison}

\begin{table}[H]
\centering
\caption{Energy Efficiency Metrics --- All Three Techniques}
\label{tab:compare_energy}
\renewcommand{\arraystretch}{1.4}
\begin{tabular}{lrrrl}
\toprule
\textbf{Metric} & \textbf{Air-Side Econ.} & \textbf{Evaporative} & \textbf{Chilled Water} & \textbf{Winner} \\
\midrule
Total Energy (kWh/yr) & 478{,}090 & \textbf{141{,}539} & 288{,}475 & Evaporative \\
Average PUE & 1.285 & \textbf{1.002} & 1.163 & Evaporative \\
Peak PUE & 1.45 & \textbf{1.004} & 1.20 & Evaporative \\
Cooling Overhead (\%) & 28.6\% & \textbf{0.2\%} & 16.3\% & Evaporative \\
Fan Energy (kWh/yr) & 105{,}965 & \textbf{366} & --- & Evaporative \\
Compressor Energy & High & \textbf{Zero} & Moderate & Evaporative \\
\bottomrule
\end{tabular}
\end{table}

\section{Environmental Impact Comparison}

\begin{table}[H]
\centering
\caption{Environmental Metrics --- All Three Techniques}
\label{tab:compare_env}
\renewcommand{\arraystretch}{1.4}
\begin{tabular}{lrrrl}
\toprule
\textbf{Metric} & \textbf{Air-Side Econ.} & \textbf{Evaporative} & \textbf{Chilled Water} & \textbf{Winner} \\
\midrule
CO$_2$ Emissions (kg/yr) & 219{,}921 & \textbf{65{,}108} & 158{,}196 & Evaporative \\
CUE (kgCO$_2$/kWh$_{IT}$) & 0.591 & \textbf{0.461} & 0.549 & Evaporative \\
Water Usage (L/yr) & \textbf{0} & \textbf{0} & 515{,}331 & Air/Evap (tie) \\
WUE (L/kWh$_{IT}$) & \textbf{0} & \textbf{0} & 1.786 & Air/Evap (tie) \\
Carbon Tax/yr (\$) & \$27{,}710 & \textbf{\$8{,}270} & \$20{,}117 & Evaporative \\
\bottomrule
\end{tabular}
\end{table}

\section{Economic Comparison}

\begin{table}[H]
\centering
\caption{Economic Metrics --- All Three Techniques}
\label{tab:compare_econ}
\renewcommand{\arraystretch}{1.4}
\begin{tabular}{lrrrl}
\toprule
\textbf{Metric} & \textbf{Air-Side Econ.} & \textbf{Evaporative} & \textbf{Chilled Water} & \textbf{Winner} \\
\midrule
CAPEX (USD) & \$45{,}000 & Low & \$630{,}000 & Evaporative \\
Annual OpEx (USD) & \$97{,}033 & \textbf{\$20{,}523} & \$42{,}553 & Evaporative \\
15-yr LCCP (USD) & --- & --- & \$1{,}838{,}582 & --- \\
NPV (USD) & \$81{,}161 & N/A & $-$\$1{,}092{,}355 & Air-Side \\
Payback (yrs) & 491.9 & N/A & 50.0 & --- \\
\bottomrule
\end{tabular}
\end{table}

\section{Thermal Compliance Comparison}

\begin{table}[H]
\centering
\caption{Thermal Compliance --- All Three Techniques}
\label{tab:compare_thermal}
\renewcommand{\arraystretch}{1.4}
\begin{tabular}{lrrrl}
\toprule
\textbf{Metric} & \textbf{Air-Side Econ.} & \textbf{Evaporative} & \textbf{Chilled Water} & \textbf{Winner} \\
\midrule
Max Inlet Temp (\degC) & $>$27 (violation) & 38.5 (violation) & \textbf{$\leq$27} & Chilled Water \\
ASHRAE A1 Compliance & \marginal & \fail & \pass & Chilled Water \\
Airflow Violations & All 5 racks & N/A & None & Chilled Water \\
Throttling Events & Yes & Yes & \textbf{None} & Chilled Water \\
Cooling Failure Hours & 0 & 0 & \textbf{0} & Tie \\
\bottomrule
\end{tabular}
\end{table}

\section{Multi-Criteria Assessment}

\begin{table}[H]
\centering
\caption{Multi-Criteria Assessment (1 = Best, 3 = Worst)}
\label{tab:radar}
\renewcommand{\arraystretch}{1.4}
\begin{tabular}{lrrr}
\toprule
\textbf{Criterion} & \textbf{Air-Side Econ.} & \textbf{Evaporative} & \textbf{Chilled Water} \\
\midrule
Energy Efficiency & 3 & \textbf{1} & 2 \\
Thermal Reliability & 2 & 3 & \textbf{1} \\
Carbon Emissions & 3 & \textbf{1} & 2 \\
Water Consumption & \textbf{1} & \textbf{1} & 3 \\
Capital Cost & 2 & \textbf{1} & 3 \\
Operating Cost & 3 & \textbf{1} & 2 \\
AI Workload Suitability & 2 & 3 & \textbf{1} \\
\midrule
\textbf{Total Score} & \textbf{16} & \textbf{11} & \textbf{14} \\
\textbf{Rank} & 3rd & \textbf{1st*} & 2nd \\
\bottomrule
\end{tabular}
\end{table}

\noindent *Evaporative ranks 1st on aggregate but fails the critical thermal compliance gate for AI workloads.

\section{Cross-Technique Comparison Graphs}

\begin{figure}[H]
\centering
\fbox{\parbox{0.75\textwidth}{\centering\vspace{2cm}
\textit{[Figure: PUE Comparison Bar Chart --- All Three Techniques]}
\vspace{2cm}}}
\caption{PUE Comparison: Air-Side Economizer vs Evaporative vs Chilled Water}
\label{fig:compare_pue}
\end{figure}

\begin{figure}[H]
\centering
\fbox{\parbox{0.75\textwidth}{\centering\vspace{2cm}
\textit{[Figure: Annual Cost Comparison --- All Three Techniques]}
\vspace{2cm}}}
\caption{Annual Operating Cost Comparison}
\label{fig:compare_cost}
\end{figure}

\begin{figure}[H]
\centering
\fbox{\parbox{0.75\textwidth}{\centering\vspace{2cm}
\textit{[Figure: Carbon Emissions Comparison --- All Three Techniques]}
\vspace{2cm}}}
\caption{Annual Carbon Emissions Comparison}
\label{fig:compare_carbon}
\end{figure}

\begin{figure}[H]
\centering
\fbox{\parbox{0.75\textwidth}{\centering\vspace{2cm}
\textit{[Figure: Water Usage Comparison --- All Three Techniques]}
\vspace{2cm}}}
\caption{Annual Water Usage Comparison}
\label{fig:compare_water}
\end{figure}

\begin{figure}[H]
\centering
\fbox{\parbox{0.75\textwidth}{\centering\vspace{2cm}
\textit{[Figure: ML Recommendation Radar --- All Three Techniques]}
\vspace{2cm}}}
\caption{ML Recommendation Composite Radar Chart}
\label{fig:compare_radar}
\end{figure}

\newpage

%% ============================================================
\chapter{Governing Equations and Formula Reference}
\label{ch:formulas}
%% ============================================================

This chapter provides a complete reference for all governing equations used across the three
simulation engines. Each formula is cross-referenced to the parameters it produces.

\section{Air-Side Economizer Equations}

\subsection{Airflow Requirement}
\begin{equation}
V_{req}\ (\text{CFM}) = \frac{P_{IT,kW} \times 3160}{\rho \cdot C_p \cdot \Delta T}
\tag{\ref{eq:air_cfm}}
\end{equation}
where $\rho = 1.2$ kg/m$^3$ (air density), $C_p = 1.006$ kJ/(kg$\cdot$K) (specific heat of air),
and $\Delta T = 15$\degC (supply-to-return temperature differential). The constant 3160 converts
kW to CFM units. \textit{Produces: \texttt{requiredAirflow\_CFM}}

\subsection{Free Cooling Heat Removal}
\begin{equation}
Q_{free} = \text{oaFraction} \times \rho \times \frac{V_{CFM}}{2118.88} \times C_p \times \Delta T_{free}
\label{eq:q_free}
\end{equation}
where $\text{oaFraction} = 1.0$ in FULL\_ECON, 0.2 in PARTIAL\_TRIM, 0.0 in MECHANICAL\_ONLY.
The constant 2118.88 converts CFM to m$^3$/s. \textit{Produces: \texttt{q\_free\_kW}}

\subsection{Mechanical Cooling Load}
\begin{equation}
Q_{mech} = \max(Q_{IT} - Q_{free},\ 0)
\label{eq:q_mech}
\end{equation}
\textit{Produces: \texttt{mech\_load\_kW}}

\subsection{Fan Power}
\begin{equation}
P_{fan} = \frac{V_{CFM} \times \eta_{fan} \times f_{filter}}{1000}
\label{eq:fan_power}
\end{equation}
where $\eta_{fan}$ is the fan efficiency coefficient and $f_{filter} = 1.15$ (15\% MERV filter penalty).
\textit{Produces: \texttt{fanPower\_kW}}

\subsection{Power Usage Effectiveness}
\begin{equation}
\text{PUE}(t) = \frac{P_{IT}(t) + P_{fan}(t) + P_{mech}(t)}{P_{IT}(t)}
\label{eq:pue}
\end{equation}
\textit{Produces: \texttt{pue} (hourly), \texttt{averagePUE} (annual)}

\subsection{Carbon Usage Effectiveness}
\begin{equation}
\text{CUE}(t) = \frac{P_{total}(t) \times \text{gridCarbonFactor}}{P_{IT}(t)}
\label{eq:cue}
\end{equation}
where gridCarbonFactor $= 0.55$ kgCO$_2$/kWh. \textit{Produces: \texttt{cue}}

\section{Evaporative Cooling Equations}

\subsection{Supply Temperature}
\begin{equation}
T_{supply} = T_{db} - \eta_{actual} \times (T_{db} - T_{wb})
\tag{\ref{eq:evap_supply}}
\end{equation}
where $\eta_{actual} \in [0.60, 0.95]$ is the saturation effectiveness.
\textit{Produces: \texttt{supplyTempC}}

\subsection{Wet-Bulb Temperature (Psychrometric)}
\begin{equation}
T_{wb} = T_{db} \times \arctan\bigl[0.151977 \times (RH + 8.313659)^{0.5}\bigr] + \ldots
\label{eq:wetbulb}
\end{equation}
(Stull 2011 approximation). \textit{Used in: Eq.~\ref{eq:evap_supply}}

\subsection{Cooling Capacity}
\begin{equation}
Q_{cooling} = \dot{V}_{air} \times \rho_{air} \times C_p \times (T_{return} - T_{supply})
\label{eq:evap_capacity}
\end{equation}
\textit{Produces: \texttt{coolingCapacityKW}}

\subsection{Thermal Mass Model}
\begin{equation}
C_{total}\frac{dT_{rack}}{dt} = Q_{IT}(t) - Q_{cooling}(t)
\tag{\ref{eq:evap_thermal_mass}}
\end{equation}
where $C_{total} = 45$ kJ/K. \textit{Produces: \texttt{inletTempC}}

\subsection{Water Evaporation Rate}
\begin{equation}
\dot{m}_{evap}\ (\text{L/hr}) = \frac{Q_{cooling} \times 3600}{L_v}
\label{eq:evap_water}
\end{equation}
where $L_v = 2{,}260$ kJ/kg. \textit{Produces: \texttt{waterEvaporationLph}}

\subsection{Thermal Throttling Model}
\begin{equation}
\Phi_{adj} = \Phi_{base} \times (1 - \gamma \cdot \Delta T)
\tag{\ref{eq:throttle}}
\end{equation}
where $\gamma = 0.10$ and $\Delta T = T_{inlet} - T_{threshold}$ (threshold = 29\degC).

\section{Chilled Water System Equations}

\subsection{EIR-Based COP}
\begin{equation}
\text{COP}_{actual} = \frac{\text{COP}_{ref}}{\text{EIR}_{FPLR}(PLR) \times \text{EIR}_{FT}(T_{chw}, T_{cond})}
\tag{\ref{eq:chw_cop}}
\end{equation}
where $\text{COP}_{ref} = 3.0$, $\text{EIR}_{FPLR}$ is a polynomial in part-load ratio (PLR),
and $\text{EIR}_{FT}$ is a biquadratic in chilled water and condenser temperatures.
\textit{Produces: \texttt{cop}}

\subsection{Chiller Power}
\begin{equation}
P_{chiller} = \frac{Q_{cooling}}{\text{COP}_{actual}}
\tag{\ref{eq:chw_power}}
\end{equation}
\textit{Produces: \texttt{chillerPower\_kW}}

\subsection{Cooling Tower Water Consumption}
\begin{equation}
\dot{V}_{tower}\ (\text{L/hr}) = \frac{Q_{rejected} \times \text{evapFraction}}{\rho_w \times C_{p,w} \times \Delta T_{range}}
\label{eq:tower_water}
\end{equation}
where $Q_{rejected} = Q_{cooling} + P_{chiller}$ (heat of rejection).
\textit{Produces: \texttt{waterUsage\_L}}

\subsection{Water Usage Effectiveness}
\begin{equation}
\text{WUE} = \frac{V_{water,annual}\ (\text{L})}{E_{IT,annual}\ (\text{kWh})}
\label{eq:wue}
\end{equation}
\textit{Produces: \texttt{wue} = 1.786 L/kWh$_{IT}$}

\subsection{Life-Cycle Cost of Plant}
\begin{equation}
\text{LCCP} = \text{CAPEX} + \sum_{y=1}^{15} \frac{\text{OpEx}_y}{(1+r)^y}
\label{eq:lccp}
\end{equation}
where $r = 0.0971$ (9.71\% discount rate). \textit{Produces: \texttt{lccp\_USD}}

\subsection{Net Present Value}
\begin{equation}
\text{NPV} = -\text{CAPEX} + \sum_{y=1}^{n} \frac{\text{Savings}_y}{(1+r)^y}
\label{eq:npv}
\end{equation}
\textit{Produces: \texttt{npv\_USD}}

\section{CloudSim Workload Model}

\subsection{Composite Workload}
\begin{equation}
U_{total}(t) = \bigl[U_{diurnal}(t) \times M_{AI}\bigr] + \sigma(t)
\tag{\ref{eq:workload}}
\end{equation}

\subsection{Diurnal Pattern}
\begin{equation}
U_{diurnal}(t) = U_{base} + A \times \sin\!\left(\frac{2\pi(t - t_{peak})}{24}\right)
\label{eq:diurnal}
\end{equation}
where $t_{peak} = 14$ (14:00 peak), $A$ is amplitude, $U_{base}$ is base utilisation.

\subsection{IT Power from Utilisation}
\begin{equation}
P_{IT}(t) = P_{idle} + (P_{peak} - P_{idle}) \times U_{total}(t)
\label{eq:it_power}
\end{equation}
\textit{Produces: \texttt{itLoad\_kW} in all three techniques}

\section{ML Scoring Function}

\subsection{Composite Score}
\begin{equation}
s_i = w_1 \cdot \hat{C}_i + w_2 \cdot \hat{E}_i + w_3 \cdot \hat{W}_i
\label{eq:ml_score}
\end{equation}
where $\hat{C}_i$, $\hat{E}_i$, $\hat{W}_i$ are min-max normalised cost, emissions, and water
for technique $i$, and $w_1 + w_2 + w_3 = 1$. Lower score = better technique.
\textit{Produces: \texttt{score} in comparison table}

\subsection{Random Forest Prediction}
\begin{equation}
\hat{y} = \arg\max_{c \in \mathcal{C}} \frac{1}{T} \sum_{t=1}^{T} \mathbf{1}\bigl[h_t(\mathbf{x}) = c\bigr]
\label{eq:rf}
\end{equation}
where $T$ is the number of trees, $h_t$ is the $t$-th decision tree, and $\mathbf{x}$ is the
feature vector. \textit{Produces: \texttt{model\_recommendation} and \texttt{model\_confidence}}

\newpage

%% ============================================================
\chapter{Discussion and Interpretation}
\label{ch:discussion}
%% ============================================================

\section{Efficiency vs Thermal Safety Trade-off}

The central finding of this study is the fundamental tension between energy efficiency and
thermal safety in data center cooling for AI workloads. Table~\ref{tab:tradeoff} summarises
this trade-off:

\begin{table}[H]
\centering
\caption{Efficiency vs Thermal Safety Trade-off Summary}
\label{tab:tradeoff}
\renewcommand{\arraystretch}{1.4}
\begin{tabular}{lcccc}
\toprule
\textbf{Technique} & \textbf{PUE} & \textbf{Thermal Safe?} & \textbf{OpEx/yr} & \textbf{Verdict} \\
\midrule
Evaporative & \textbf{1.002} & \textcolor{failred}{No (38.5\degC)} & \textbf{\$20{,}523} & Efficient but unsafe \\
Chilled Water & 1.163 & \textcolor{passgreen}{Yes ($\leq$27\degC)} & \$42{,}553 & Safe but costly \\
Air-Side Econ. & 1.285 & \textcolor{warnorng}{Marginal} & \$97{,}033 & Climate-dependent \\
\bottomrule
\end{tabular}
\end{table}

\section{Implications for AI Workload Data Centers}

AI training workloads impose sustained high-density thermal loads that fundamentally differ
from traditional enterprise IT. The simulation results demonstrate three key implications:

\begin{enumerate}
  \item \textbf{Airflow-based cooling reaches physical limits:} At AI workload densities
  ($>$10 kW/rack), the required airflow (CFM) exceeds the physical envelope of air-side
  economization. The simulation recorded violations at all five racks during peak hours.
  
  \item \textbf{Evaporative cooling is a pre-cooling candidate, not a standalone solution:}
  The 38.5\degC inlet temperature (11.5\degC above ASHRAE A1 limit) makes standalone
  evaporative cooling unsuitable. However, its near-zero overhead (PUE 1.002) makes it
  an excellent pre-cooling stage in a hybrid configuration.
  
  \item \textbf{Chilled water is the only thermally compliant solution:} The chilled water
  system is the only technique that maintains ASHRAE Class A2 compliance throughout the
  8{,}760-hour simulation. Its higher CAPEX (\$630{,}000) and water consumption (515{,}331 L/yr)
  are the trade-offs for this reliability.
\end{enumerate}

\section{ML Recommendation vs Engineering Judgment}

The ML model recommends Evaporative Cooling (94.75\% confidence) based on cost and emissions
metrics. However, engineering judgment requires overriding this recommendation for AI-density
deployments due to the thermal compliance failure. This highlights a key limitation of the
current ML pipeline: the scoring function does not include a hard thermal compliance constraint.

\begin{table}[H]
\centering
\caption{ML Recommendation vs Engineering Judgment}
\label{tab:ml_vs_eng}
\renewcommand{\arraystretch}{1.3}
\begin{tabular}{lll}
\toprule
\textbf{Criterion} & \textbf{ML Recommendation} & \textbf{Engineering Judgment} \\
\midrule
Primary metric & Cost + Emissions + Water & Thermal compliance (hard constraint) \\
Recommended technique & Evaporative & Chilled Water \\
Confidence & 94.75\% & N/A (deterministic) \\
Thermal compliance & Not penalised & Hard gate (FAIL = disqualify) \\
Suitable for AI loads & No (thermal failure) & Yes (chilled water only) \\
\bottomrule
\end{tabular}
\end{table}

\section{Hybrid Configuration Recommendation}

Based on the simulation results, the optimal configuration for AI-density data centers is a
\textbf{hybrid evaporative + chilled water system}:

\begin{enumerate}[noitemsep]
  \item Deploy evaporative cooling as a pre-cooling stage to reduce condenser inlet temperature.
  \item Use chilled water as the primary cooling loop to guarantee ASHRAE compliance.
  \item Activate air-side economization during favorable weather windows (T$_{out} \leq 18$\degC).
  \item Expected hybrid PUE: 1.05--1.10 (between evaporative and chilled water standalone values).
\end{enumerate}

\newpage

%% ============================================================
\chapter{Conclusion}
\label{ch:conclusion}
%% ============================================================

This chapter has presented the complete simulation results for three data center cooling
strategies evaluated over 8{,}760 hours under AI-era workload conditions. The key conclusions are:

\begin{enumerate}
  \item \textbf{Evaporative Cooling} achieves the best energy efficiency (PUE 1.002) and lowest
  operating cost (\$20{,}523/yr) but fails ASHRAE thermal compliance (inlet temp 38.5\degC).
  It is recommended as a pre-cooling stage in hybrid configurations.
  
  \item \textbf{Chilled Water System} is the only technique achieving full ASHRAE Class A2
  thermal compliance. Its high CAPEX (\$630{,}000) and water consumption (515{,}331 L/yr)
  are justified for high-density AI workloads where thermal reliability is mandatory.
  
  \item \textbf{Air-Side Economization} provides 28.6\% energy savings and zero water consumption
  but is physically limited by airflow constraints at AI workload densities. It is suitable
  for moderate-density deployments in cool, dry climates.
  
  \item \textbf{The ML recommender} (Random Forest, 94.75\% confidence) recommends Evaporative
  Cooling based on cost and emissions metrics. Engineering judgment overrides this for AI-density
  deployments due to the thermal compliance failure.
  
  \item \textbf{The optimal strategy} for AI data centers is a hybrid evaporative + chilled water
  configuration, combining the efficiency advantages of evaporative pre-cooling with the thermal
  reliability of vapour-compression refrigeration.
\end{enumerate}

\newpage

%% ============================================================
\appendix
\chapter{Raw Simulation Data Samples}
\label{app:raw}
%% ============================================================

\section{Air-Side Economizer --- Extended Hourly Sample (Hours 0--23)}

\begin{longtable}{rrrrrrrr}
\caption{Air-Side Economizer --- Full 24-Hour Hourly Results} \label{tab:air_24hr} \\
\toprule
\textbf{Hr} & \textbf{T$_{out}$} & \textbf{RH} & \textbf{IT (kW)} & \textbf{CFM} & \textbf{Mode} & \textbf{Total (kW)} & \textbf{PUE} \\
\midrule
\endfirsthead
\multicolumn{8}{c}{\tablename\ \thetable{} -- continued} \\
\toprule
\textbf{Hr} & \textbf{T$_{out}$} & \textbf{RH} & \textbf{IT (kW)} & \textbf{CFM} & \textbf{Mode} & \textbf{Total (kW)} & \textbf{PUE} \\
\midrule
\endhead
\midrule \multicolumn{8}{r}{Continued on next page} \\
\endfoot
\bottomrule
\endlastfoot
0 & 24.00 & 60.00 & 35.00 & 7634 & FULL\_ECON & 43.23 & 1.235 \\
1 & 24.01 & 59.99 & 34.13 & 7446 & FULL\_ECON & 42.18 & 1.236 \\
2 & 24.01 & 59.98 & 33.84 & 7382 & FULL\_ECON & 41.82 & 1.236 \\
3 & 24.02 & 59.97 & 34.13 & 7446 & FULL\_ECON & 42.19 & 1.236 \\
4 & 24.02 & 59.96 & 35.00 & 7634 & FULL\_ECON & 43.26 & 1.236 \\
5 & 24.03 & 59.95 & 36.37 & 7934 & FULL\_ECON & 44.97 & 1.237 \\
6 & 24.03 & 59.94 & 38.16 & 8324 & FULL\_ECON & 47.19 & 1.237 \\
7 & 24.04 & 59.92 & 40.24 & 8779 & FULL\_ECON & 49.78 & 1.237 \\
8 & 24.05 & 59.91 & 42.48 & 9266 & FULL\_ECON & 52.56 & 1.237 \\
9 & 24.05 & 59.90 & 44.72 & 9754 & FULL\_ECON & 55.33 & 1.237 \\
10 & 24.06 & 59.89 & 46.80 & 10209 & FULL\_ECON & 57.93 & 1.238 \\
11 & 24.06 & 59.88 & 48.59 & 10599 & FULL\_ECON & 60.14 & 1.238 \\
12 & 24.07 & 59.87 & 49.96 & 10899 & FULL\_ECON & 61.86 & 1.238 \\
13 & 24.07 & 59.86 & 50.83 & 11087 & FULL\_ECON & 62.93 & 1.238 \\
14 & 24.08 & 59.85 & 51.12 & 11151 & FULL\_ECON & 63.32 & 1.239 \\
15 & 24.09 & 59.84 & 50.83 & 11087 & FULL\_ECON & 62.97 & 1.239 \\
16 & 24.09 & 59.83 & 49.96 & 10899 & FULL\_ECON & 61.90 & 1.239 \\
17 & 24.10 & 59.82 & 48.59 & 10599 & FULL\_ECON & 60.22 & 1.239 \\
18 & 24.10 & 59.81 & 46.80 & 10209 & FULL\_ECON & 58.01 & 1.239 \\
19 & 24.11 & 59.80 & 44.72 & 9754 & FULL\_ECON & 55.44 & 1.240 \\
20 & 24.11 & 59.78 & 42.48 & 9266 & FULL\_ECON & 52.67 & 1.240 \\
21 & 24.12 & 59.77 & 40.24 & 8779 & FULL\_ECON & 49.91 & 1.240 \\
22 & 24.13 & 59.76 & 38.16 & 8324 & FULL\_ECON & 47.34 & 1.241 \\
23 & 24.13 & 59.75 & 36.37 & 7934 & FULL\_ECON & 45.12 & 1.241 \\
\end{longtable}

\noindent\textbf{Note:} Hours 10--18 (shaded) exceed the 10{,}000 CFM airflow limit, triggering
violation flags and engineering recommendations to upgrade to liquid cooling.

\section{Evaporative Cooling --- Extended Hourly Sample (Hours 0--9)}

\begin{table}[H]
\centering
\caption{Evaporative Cooling --- Extended Hourly Results (Hours 0--9)}
\label{tab:evap_10hr}
\renewcommand{\arraystretch}{1.2}
\scriptsize
\begin{tabular}{rrrrrrrrr}
\toprule
\textbf{Hr} & \textbf{T$_{amb}$} & \textbf{IT (kW)} & \textbf{Fan (kW)} & \textbf{Total (kW)} & \textbf{PUE} & \textbf{T$_{inlet}$} & \textbf{T$_{supply}$} & \textbf{Mode} \\
\midrule
0 & 25.0 & 11.760 & 0.01424 & 11.774 & 1.0012 & 33.874 & 18.874 & DEC \\
1 & 25.0 & 11.760 & 0.01424 & 11.774 & 1.0012 & 33.882 & 18.882 & DEC \\
2 & 25.0 & 11.760 & 0.01424 & 11.774 & 1.0012 & 33.880 & 18.880 & DEC \\
3 & 25.0 & 11.760 & 0.01424 & 11.774 & 1.0012 & 33.888 & 18.888 & DEC \\
4 & 25.0 & 11.760 & 0.01424 & 11.774 & 1.0012 & 33.874 & 18.874 & DEC \\
5 & 25.0 & 11.798 & 0.01444 & 11.813 & 1.0012 & 33.882 & 18.882 & DEC \\
6 & 25.0 & 12.390 & 0.01516 & 12.405 & 1.0012 & 33.888 & 18.888 & DEC \\
7 & 25.0 & 13.034 & 0.01595 & 13.050 & 1.0012 & 33.895 & 18.895 & DEC \\
8 & 25.0 & 13.720 & 0.01679 & 13.737 & 1.0012 & 33.901 & 18.901 & DEC \\
9 & 25.0 & 14.440 & 0.01767 & 14.458 & 1.0012 & 33.908 & 18.908 & DEC \\
\bottomrule
\end{tabular}
\end{table}

\section{Chilled Water System --- Extended Hourly Sample (Hours 0--9)}

\begin{table}[H]
\centering
\caption{Chilled Water System --- Extended Hourly Results (Hours 0--9)}
\label{tab:chw_10hr}
\renewcommand{\arraystretch}{1.2}
\scriptsize
\begin{tabular}{rrrrrrrrr}
\toprule
\textbf{Hr} & \textbf{T$_{amb}$} & \textbf{IT (kW)} & \textbf{Cool (kW)} & \textbf{Chiller (kW)} & \textbf{COP} & \textbf{Water (L)} & \textbf{Cost (\$)} & \textbf{CO$_2$ (kg)} \\
\midrule
0 & 24.00 & 21.426 & 2.757 & 2.678 & 8.000 & 43.387 & 2.636 & 14.122 \\
1 & 24.01 & 21.426 & 2.757 & 2.678 & 7.999 & 43.387 & 2.636 & 14.122 \\
2 & 24.01 & 21.426 & 2.757 & 2.679 & 7.999 & 43.388 & 2.636 & 14.122 \\
3 & 24.02 & 21.426 & 2.758 & 2.679 & 7.998 & 43.388 & 2.636 & 14.122 \\
4 & 24.02 & 21.426 & 2.758 & 2.679 & 7.998 & 43.389 & 2.636 & 14.123 \\
5 & 24.03 & 24.622 & 3.198 & 3.079 & 7.997 & 49.862 & 2.942 & 15.759 \\
6 & 24.03 & 25.092 & 3.265 & 3.138 & 7.996 & 50.815 & 4.267 & 16.001 \\
7 & 24.04 & 25.806 & 3.365 & 3.228 & 7.996 & 52.260 & 4.364 & 16.367 \\
8 & 24.05 & 26.745 & 3.498 & 3.345 & 7.995 & 54.162 & 4.493 & 16.849 \\
9 & 24.05 & 27.793 & 3.649 & 3.477 & 7.994 & 56.285 & 4.637 & 17.389 \\
\bottomrule
\end{tabular}
\end{table}

\newpage

%% ============================================================
\chapter{Figure Reference Summary}
\label{app:figures}
%% ============================================================

This appendix lists all figure placeholders in this document. Replace each placeholder with
the corresponding screenshot or export from the frontend dashboard.

\begin{table}[H]
\centering
\caption{Complete Figure Reference List}
\label{tab:fig_ref}
\renewcommand{\arraystretch}{1.3}
\begin{tabular}{p{2cm}p{5cm}p{7cm}}
\toprule
\textbf{Figure} & \textbf{Caption} & \textbf{Source / Frontend Location} \\
\midrule
\ref{fig:cosim_arch} & Co-Simulation Architecture & Architecture diagram (draw manually or export from docs) \\
\ref{fig:air_power} & Air-Side: Hourly Power Breakdown & Air-Side Economization tab $\to$ Power Breakdown chart \\
\ref{fig:air_airflow} & Air-Side: Airflow vs Cooling Split & Air-Side Economization tab $\to$ Airflow chart \\
\ref{fig:air_modes} & Air-Side: Mode Distribution & Air-Side Economization tab $\to$ Mode Distribution chart \\
\ref{fig:air_cost} & Air-Side: Cost Structure & Air-Side Economization tab $\to$ Cost Structure chart \\
\ref{fig:air_racks} & Air-Side: Rack Analysis & Air-Side Economization tab $\to$ CloudSim Rack Analysis \\
\ref{fig:air_pue_cue} & Air-Side: PUE and CUE & Air-Side Economization tab $\to$ PUE/CUE chart \\
\ref{fig:air_ambient} & Air-Side: Ambient Conditions & Air-Side Economization tab $\to$ Ambient Conditions chart \\
\ref{fig:air_radar} & Air-Side: Performance Radar & Air-Side Economization tab $\to$ Radar chart \\
\ref{fig:air_5yr} & Air-Side: 5-Year Projection & Air-Side Economization tab $\to$ 5-Year Projection chart \\
\ref{fig:evap_summary} & Evap: Annual Summary & Evaporative Cooling tab $\to$ Annual Summary \\
\ref{fig:evap_power} & Evap: Hourly Power Breakdown & Evaporative Cooling tab $\to$ Power Breakdown chart \\
\ref{fig:evap_capacity} & Evap: Capacity vs IT Load & Evaporative Cooling tab $\to$ Capacity vs Load chart \\
\ref{fig:evap_pue} & Evap: PUE vs Annual Average & Evaporative Cooling tab $\to$ PUE chart \\
\ref{fig:evap_temp_hum} & Evap: Temperature and Humidity & Evaporative Cooling tab $\to$ Temp/Humidity chart \\
\ref{fig:evap_supply} & Evap: Supply Air Conditions & Evaporative Cooling tab $\to$ Supply Air chart \\
\ref{fig:evap_modes} & Evap: Mode Distribution & Evaporative Cooling tab $\to$ Mode Distribution \\
\ref{fig:evap_pue_hourly} & Evap: Hourly PUE Trace & Evaporative Cooling tab $\to$ Hourly PUE chart \\
\ref{fig:evap_assess_chart} & Evap: Cooling Assessment & Evaporative Cooling tab $\to$ Assessment checklist \\
\ref{fig:evap_radar} & Evap: Performance Radar & Evaporative Cooling tab $\to$ Radar chart \\
\ref{fig:evap_5yr} & Evap: 5-Year Projection & Evaporative Cooling tab $\to$ 5-Year Projection \\
\ref{fig:chw_overview} & CHW: Annual Overview & Chilled Water tab $\to$ Annual Overview \\
\ref{fig:chw_cop} & CHW: COP Over Time & Chilled Water tab $\to$ COP chart \\
\ref{fig:chw_load_power} & CHW: Cooling Load vs Chiller Power & Chilled Water tab $\to$ Load vs Power chart \\
\ref{fig:chw_it_chiller} & CHW: IT Load vs Chiller Power & Chilled Water tab $\to$ IT vs Chiller chart \\
\ref{fig:chw_water_carbon} & CHW: Water and Carbon Trend & Chilled Water tab $\to$ Water/Carbon chart \\
\ref{fig:chw_cost} & CHW: Cost Structure & Chilled Water tab $\to$ Cost Structure chart \\
\ref{fig:chw_gates} & CHW: Phase 4 Gates & Chilled Water tab $\to$ Compliance Gates \\
\ref{fig:chw_radar} & CHW: Performance Radar & Chilled Water tab $\to$ Radar chart \\
\ref{fig:chw_5yr} & CHW: 5-Year Projection & Chilled Water tab $\to$ 5-Year Projection \\
\ref{fig:compare_pue} & Comparison: PUE & Results Comparison tab $\to$ PUE bar chart \\
\ref{fig:compare_cost} & Comparison: Annual Cost & Results Comparison tab $\to$ Cost bar chart \\
\ref{fig:compare_carbon} & Comparison: Carbon Emissions & Results Comparison tab $\to$ Emissions chart \\
\ref{fig:compare_water} & Comparison: Water Usage & Results Comparison tab $\to$ Water chart \\
\ref{fig:compare_radar} & Comparison: ML Radar & Results Comparison tab $\to$ ML Radar chart \\
\bottomrule
\end{tabular}
\end{table}

\end{document}
